\chapter*{平面直角坐标系中的直线}

\section*{直线的倾斜角和斜率}

\zsd 倾斜角、斜率

A 组：

\begin{enumerate}
    \item 已知直线过点(6,0)和 $\left( {3,\sqrt{3}}\right)$ ,则直线的倾斜角为\blkda{$\ang{150}$}.
    
    \item 若直线的倾斜角为 ${2a} - \dfrac{\pi }{6}$ ,则 $a$ 的取值范围是\blkda{$\left[\dfrac{\pi}{12}, \dfrac{7\pi}{12}\right)$}.
    \begin{jieda}
        $2a-\dfrac{\pi}{6} \in [0, \pi)$，故$a \in \left[\dfrac{\pi}{12}, \dfrac{7\pi}{12}\right)$
    \end{jieda}
    
    \item 已知直线 $l$ 平行于第二、四象限的角平分线,则直线 $l$ 的倾斜角为\blkda{$\dfrac{3\pi}{4}$}(用弧度制表示).
    
    \item 经过 $A\left( {-1,3}\right) ,B\left( {\sqrt{3}, - \sqrt{3}}\right)$ 两点的直线的斜率 $k =$ \blkda{$-\sqrt{3}$},其倾斜角 $\alpha =$ \blkda{$\dfrac{2\pi}{3}$}.
    \begin{jieda}
        $k = \dfrac{-\sqrt{3}-3}{\sqrt{3} - (-1)} = -\sqrt{3}$，$tan\alpha = k = -\sqrt{3}$，$\alpha = \dfrac{2\pi}{3}$.
    \end{jieda}
    
    \item 已知直线 $l$ 过 $A\left( {a,2}\right)$ 、 $B\left( {2, - 4}\right)$ 两点，求直线 $l$ 的倾斜角 $\alpha$.
    \begin{jieda}
        \begin{dingautolist}{172}
            \item $a = 2$时，$\alpha = \dfrac{\pi}{2}$
            \item $a > 2$时，$\alpha = arctan\dfrac{6}{a-2}$
            \item $a < 2$时，$\alpha = \pi - arctan\dfrac{6}{2-a}$
        \end{dingautolist}
    \end{jieda}
    
    \item 直线 $l : x = \tan \dfrac{\pi }{4}$ 的倾斜角 $a =$ \blkda{$\dfrac{\pi}{2}$},斜率 $k =$ \blkda{不存在}.
    
    \item 若直线 $l$ 的倾斜角 $a \in  \left\lbrack  {\dfrac{\pi }{3},\dfrac{\pi }{2}}\right)  \cup  \left( {\dfrac{\pi }{2},\dfrac{5\pi }{6}}\right\rbrack$ ,则它的斜率 $k$ 的取值范围是\blkda[4]{$\left(-\infty, -\dfrac{\sqrt{3}}{3}\right] \bigcup \left[\sqrt{3}, +\infty\right)$}. 
\end{enumerate}

B 组：
\begin{enumerate}
    \item 已知过点 $A\left( {1,a}\right)$ , $B\left( {-2,\sqrt{3}}\right)$ 的直线的倾斜角为 ${60}^{ \circ  }$ ,则实数 $a =$ \blkda{$4\sqrt{3}$}.
    \begin{jieda}
        直线$AB$的斜率的$k = \dfrac{\sqrt{3}-a}{-2-1} = tan\ang{60} = \sqrt{3}$，解得$a = 4\sqrt{3}$.
    \end{jieda}
    
    \item 直线 $l : \dfrac{x - 2}{3} = \dfrac{y + 3}{-2}$ 的斜率是\blkda{$-\dfrac{2}{3}$},其倾斜角 $\alpha =$ \blkda[2]{$\pi - arctan\dfrac{2}{3}$}.
    \begin{jieda}
        $\dfrac{x - 2}{3} = \dfrac{y + 3}{-2} \Longleftrightarrow 4-2x = 3y+9 \Longleftrightarrow y = -\dfrac{2}{3}x - \dfrac{5}{3}$，故其斜率为$-\dfrac{2}{3}$，倾斜角为$\pi - arctan\dfrac{2}{3}$
    \end{jieda}

    \item 过 $A\left( {-3,2}\right)$ 、 $B\left( {2,1}\right)$ 两点的直线的倾斜角为\blkda[2]{$\pi - arctan\dfrac{1}{5}$}.

    \item 若 $\vv{n} = \left( {a,b}\right) \left( {{ab} > 0}\right)$ 是与直线 $l$ 垂直方向的一个向量,则直线 $l$ 的斜率 $k =$ \blkda{$-\dfrac{a}{b}$},倾斜角 $a =$ \blkda{$\pi - arctan\dfrac{a}{b}$}.
    \begin{jieda}
        $ax+by+c = 0$
    \end{jieda}

    \item 已知直线 $l$ 过点 $P\left( {-1,2}\right)$ ,且它与以 $A\left( {-2, - 3}\right) \text{ 、 }B\left( {3,0}\right)$ 为端点的线段 ${AB}$ 有公共 点, 则直线 $l$ 的斜率的取值范围是\blkda[3]{$\left( - \infty, -\dfrac{1}{2}\right\rbrack   \cup  \lbrack 5, + \infty )$}.
    \begin{jieda}
        $\because$ 点 $P\left( {-1,2}\right)$ 、 $A\left( {-2, - 3}\right)$
        
        $\therefore$ 直线 ${AP}$ 的斜率 ${k}_{1} = \dfrac{-3 - 2}{-2 + 1} = 5$ .同理可得直线 ${BP}$ 的斜率 ${k}_{2} =  -\dfrac{1}{2}$ .
        \begin{center}
            \includegraphics[width=0.35\linewidth]{bodyxb//src/zhixiandeqingjiao1.png}
        \end{center}
        设直线 $l$ 与线段 ${AB}$ 交于$M$ 点，当直线的倾斜角为锐角时，随着 $M$ 从 $A$ 向 $B$ 移动的过程中, $l$ 的倾斜角变大，$l$ 的斜率也变大,直到 ${PM}$ 平行 $y$ 轴时 $l$ 的斜率不存在，此时 $l$ 的斜率 $k \geq  5$ ；
    
        当直线的倾斜角为钝角时,随着 $l$ 的倾斜角变大, $l$ 的斜率从负无穷增大到直线 ${BP}$ 的斜率,此时 $l$ 的斜率 $k \leq   - \dfrac{1}{2}$ .
    
        综上所述,可得直线 $l$ 的斜率取值范围为： $\left( - \infty, -\dfrac{1}{2}\right\rbrack   \cup  \lbrack 5, + \infty )$ .
    \end{jieda}

    \item 已知 $A\left( {{\tan }^{2}a,0}\right)$ 、 $B\left( {1,2\tan a}\right)$ 两点,且 $0 < a < \dfrac{\pi }{4}$ . 直线 ${AB}$ 的倾斜角(用 $a$ 表示)为\blkda{$2a$}.
    \begin{jieda}
        直线$AB$的斜率$k = \dfrac{2tana}{1-tan^2a} = tan2a$，因为$0 < a < \dfrac{\pi }{4}$，所以直线 ${AB}$ 的倾斜角为$2a$.
    \end{jieda}

    \item 经过 $A\left( {1,m}\right) ,B\left( {m - 1,3}\right)$ 两点的直线的倾斜角是钝角,则实数 $m$ 的范围是\blkda[3]{$(-\infty, 2) \cup (3, +\infty)$}.
    \begin{jieda}
        $\vv{AB} = (m-2, 3-m)$，当$(m-2)(3-m) < 0$时满足要求，故$m \in (-\infty, 2) \cup (3, +\infty)$
    \end{jieda}
    
    \item 直线 ${ax} + {by} + c = 0\left( {a > 0,b > 0}\right)$ 的倾斜角是\blkda[3]{$\pi - arctan\dfrac{a}{b}$}.
    
    \item 直线 $x\cos \alpha  - y + 1 = 0$ 的倾斜角的范围是\blkda[4]{$\left[0, \dfrac{\pi}{4}\right] \cup \left[\dfrac{3\pi}{4}, \pi \right)$}.
    \begin{jieda}
        $k = cos\alpha \in [-1, 1]$
    \end{jieda}

    \item 若点 $M\left( {x,y}\right)$ 在函数 $y =  - {2x} + 8$ 的图像上,当 $x \in  \left\lbrack  {2,5}\right\rbrack$ 时,则 $\dfrac{y + 1}{x + 1}$ 的取值范围是\blkda{$\left[-\dfrac{1}{6}, \dfrac{5}{3}\right]$}.
    \begin{jieda}
        $\dfrac{y+1}{x+1} = \dfrac{9-2x}{x+1} = \dfrac{11}{x+1}-2 \in \left[-\dfrac{1}{6}, \dfrac{5}{3}\right]$
    
        \ps 可以看做直线上的点到点$(-1, -1)$连线的斜率取值范围
    \end{jieda}

    \item 在平面直角坐标系中, 横坐标、纵坐标都为有理数的点称为有理点. 如果任作一条倾斜角为 $\dfrac{\pi }{3}$ 的直线,那么这样的直线最多能通过\blkda{1}个有理点.

    \item 若直线 $l$ 经过点 $P\left( {0,m}\right)$ 且它的倾斜角为 ${60}^{ \circ  }$ ,则直线 $l$ 绕点 $P$ 沿顺时针方向转过 ${90}^{ \circ  }$ 后所得直线的斜率为\blkda{$-\dfrac{\sqrt{3}}{3}$}.

    \item 下列结论中,正确的是\dotfill{(\kh{D})}
    \xx{一条直线的倾斜角的正切叫做直线的斜率}{若一条直线经过 $\left( {{x}_{1},{y}_{1}}\right) ,\left( {{x}_{2},{y}_{2}}\right)$ 两点,则这条直线的斜率 $k = \dfrac{{y}_{2} - {y}_{1}}{{x}_{2} - {x}_{1}}$}{直线的倾斜角的变化范围是 $\left\lbrack  {0,\dfrac{\pi }{2}}\right)  \cup  \left( {\dfrac{\pi }{2},\pi }\right)$}{若直线 $l$ 平行于 $x$ 轴,则直线 $l$ 的倾斜角为零}

    \item 图中的直线 ${l}_{1}\text{ 、 }{l}_{2}\text{ 、 }{l}_{3}$ 的斜率分别为 ${k}_{1}\text{ 、 }{k}_{2}\text{ 、 }{k}_{3}$ ,则\dotfill{(\kh{D})}
    
    \begin{center}
    \includegraphics[width=0.3\textwidth]{bodyxb/src/bo_d2nqu4ref24c73b1hs1g_1_240_1400_333_286_0.jpg}
    \end{center}

    \xx{${k}_{1} < {k}_{2} < {k}_{3}$}{${k}_{3} < {k}_{1} < {k}_{2}$}{${k}_{3} < {k}_{2} < {k}_{1}$}{${k}_{1} < {k}_{3} < {k}_{2}$}
    \begin{jieda}
        设直线 ${l}_{1},{l}_{2},{l}_{3}$ 的倾斜角分别为 ${a}_{1},{a}_{2},{a}_{3}$ 
        
        由图像可得 ${a}_{1} > {90}^{ \circ  } > {a}_{2} > {a}_{3} > {0}^{ \circ  }$ ,由倾斜角与斜率的关系可得${k}_{1} < 0 < {k}_{3} < {k}_{2}$ .
    \end{jieda}

    \item 若直线 $l$ 的斜率 $k$ 满足 $- 2 \leq  k \leq  2$ ,则其倾斜角 $a$ 的取值范围是\dotfill{(\kh{C})}
    \xx{$- \arctan 2 \leq  a \leq  \arctan 2$}{$\arctan 2 \leq  a \leq  \pi  - \arctan 2$}{$0 \leq  a \leq  \arctan 2$ 或 $\pi  - \arctan 2 \leq  a < \pi$}{$0 \leq  a < \arctan 2$ 或 $\pi  - \arctan 2 < a < \pi$}

    \item 若 $- \dfrac{\pi }{2} < a < 0$ ,则直线 $\dfrac{x}{\sin a} = \dfrac{y}{\cos a}$ 的倾斜角为\dotfill{(\kh{A})}
    \xx{$\dfrac{\pi }{2} - a$}{$a - \dfrac{\pi }{2}$}{$\pi  - a$}{$\dfrac{3\pi }{2} - a$}
    \begin{jieda}
        $k = cota < 0$
    \end{jieda}
    
    \item 下面真命题的个数是\blkda{0}个.
    \begin{dingautolist}{172}
        \item 在平面直角坐标系内,任何一条直线都有倾斜角和斜率；
        \item 若一条直线倾斜角是另一条的两倍,则斜率也是他的两倍；
        \item 直线 $y = 1$ 的斜率不存在;
        \item 经过 ${A}\left( {{x}_{1},{y}_{1}}\right) ,{B}\left( {{x}_{2},{y}_{2}}\right)$ 的直线斜率是 $k = \dfrac{{y}_{2} - {y}_{1}}{{x}_{2} - {x}_{1}}$
    \end{dingautolist}
    
    \item 若 $a \in  \left\lbrack  {\dfrac{\pi }{6},\dfrac{\pi }{2}}\right)$ ,则直线 ${2x}\cos a + {3y} + 1 = 0$ 的倾斜角的取值范围是\dotfill{(\kh{B})}
    \xx{$\left\lbrack  {\dfrac{\pi }{6},\dfrac{\pi }{2}}\right)$}{$\left\lbrack  {\dfrac{5\pi }{6},\pi }\right)$}{$\left\lbrack  {0,\dfrac{\pi }{6}}\right)$}{$\left\lbrack  {\dfrac{\pi }{2},\dfrac{5\pi }{6}}\right)$}
    \begin{jieda}
        $a \in  \left\lbrack  {\dfrac{\pi }{6},\dfrac{\pi }{2}}\right)$，故$cosa \in \left(0, \dfrac{\sqrt{3}}{2}\right]$
        
        $k = -\dfrac{2cosa}{3} \in \left[-\dfrac{\sqrt{3}}{3}, 0\right)$，故倾斜角的取值范围是$\left[\dfrac{5\pi}{6}, \pi \right)$
    \end{jieda}
    
    \item 已知直线 $l$ 的倾斜角的正弦是 $\dfrac{3}{5}$ ,且直线 $l$ 过 $A\left( {x, - 3}\right) \text{ 、 }B\left( {2,y}\right) \text{ 、 }C\left( {5,3}\right)$ 三点, 求 $x$ 、 $y$ 的值.
    \begin{jieda}
        $\left\{  {\begin{array}{l} x =  - 3 \\  y = \dfrac{3}{4} \end{array}\text{ 或 }\left\{  \begin{array}{l} x = {13} \\  y = \dfrac{21}{4} \end{array}\right. }\right.$
    \end{jieda}

    \item 已知两个不同的点 $A\left( {-\cos a,3{\sin }^{2}a}\right)$ 、 $B\left( {\sqrt{3}\sin a,{\cos }^{2}a}\right)$ ,且 $a \in  \left\lbrack  {\dfrac{\pi }{2},\pi }\right\rbrack$ ,求直线 ${AB}$ 的斜率及倾斜角的取值范围.
    \begin{jieda}
        $k_{AB} = \dfrac{cos^2a - 3sin^2a}{\sqrt{3}sina + cosa} =  (cosa - \sqrt{3}sina) \times \dfrac{2sin\left(a+\dfrac{\pi}{6}\right)}{2sin\left(a+\dfrac{\pi}{6}\right)} = -2sin\left(a - \dfrac{\pi}{6}\right), a \in \left[\dfrac{\pi}{2}, \dfrac{5\pi}{6}\right)\cup \left(\dfrac{5\pi}{6}, \pi\right]$（检验可知当$a = \dfrac{5\pi}{6}$时两点重合不符合题意）

        故$k_{AB} \in \left[-2, -1\right]$，因此倾斜角的取值范围是$\left[\pi - arctan2, \dfrac{3\pi}{4}\right]$
    \end{jieda}

\end{enumerate}
C 组：
\begin{enumerate}
    \item 若点 $M\left( {x,y}\right)$ 在一次函数 $y =  - {2x} + 8$ 的图像上,当 $x \in  \left\lbrack  {2,5}\right\rbrack$ 时,则 $\dfrac{{2y} + 1}{x + 1}$ 的取值范围是\blkda{$\left\lbrack  {-\dfrac{1}{2},3}\right\rbrack$}.
    \begin{jieda}
        如图,
        
        \begin{center}
        \includegraphics[width=0.3\textwidth]{bodyxb/src/bo_d2ntcqref24c73b1ilq0_2_238_638_757_804_0.jpg}
        \end{center}
        
        函数 $y =  - {2x} + 8,x \in  \left\lbrack  {2,5}\right\rbrack$ 表示线段 ${AB}$ 其中 $A\left( {5, - 2}\right) ,B\left( {2,4}\right)$ ,
        
        $\dfrac{{2y} + 1}{x + 1} = \dfrac{2\left( {y + \dfrac{1}{2}}\right) }{x + 1}$ 的几何意义为线段 ${AB}$ 上的动点 $M\left( {x,y}\right)$ 与定点 $P\left( {-1, - \dfrac{1}{2}}\right)$ 连线的斜率的 2 倍,
        
        $\because {k}_{PA} = \dfrac{-2 + \dfrac{1}{2}}{5 + 1} =  - \dfrac{1}{4},{k}_{PB} = \dfrac{4 + \dfrac{1}{2}}{2 + 1} = \dfrac{3}{2}$ ,
        
        $\therefore  - \dfrac{1}{4} \leq  {k}_{PM} \leq  \dfrac{3}{2}$
        
        $\therefore \dfrac{{2y} + 1}{x + 1}$ 的取值范围是 $\left\lbrack  {-\dfrac{1}{2},3}\right\rbrack$
    \end{jieda}
    \item 已知 $k$ 为正实数,设直线 ${l}_{1}$ 的斜率为 $k$ ,直线 ${l}_{2}$ 的斜率为 $\dfrac{2}{k}$ ,且 ${l}_{1}$ 与 ${l}_{2}$ 交于 $x$ 轴外一点,若 ${l}_{1},{l}_{2}$ 与 $x$ 轴围成一个等腰三角形,则 $k$ 的所有可能的取值集合为\blkda{$\left\{  {2\sqrt{2},\dfrac{\sqrt{2}}{2}}\right\}$}.
    \begin{jieda}
        设直线 ${l}_{1}$ 与直线 ${l}_{2}$ 的倾斜角为 $\alpha ,\beta$ ,
        
        因为 $k$ 为正实数,所以 $\alpha ,\beta$ 均为锐角,
        
        因为直线 ${l}_{1},{l}_{2}$ 与 $x$ 轴围成一个等腰三角形,则有以下两种情况,
        
        (1) $a = {2\beta }$ 时, $\tan \alpha  = \tan {2\beta } = \dfrac{2\tan \beta }{1 - {\tan }^{2}\beta }$ ,所以 $k = \dfrac{\dfrac{4}{k}}{1 - {\left( \dfrac{2}{k}\right) }^{2}}$ ,
        
        因为 $k > 0$ ,解得 $k = 2\sqrt{2}$ ,
        
        (2) ${2a} = \beta$ 时, $\tan \beta  = \tan {2a} = \dfrac{2\tan a}{1 - {\tan }^{2}a}$ ,所以 $\dfrac{2}{k} = \dfrac{2k}{1 - {k}^{2}}$ ,
        
        因为 $k > 0$ ,解得 $k = \dfrac{\sqrt{2}}{2}$ .
    \end{jieda}
    \item 1949 年公布的《国旗制法说明》中就五星的位置规定:大五角星有一个角尖正向上方, 四颗小五角星均各有一个角尖正对大五角星的中心点. 有人发现, 第三颗小星的姿态与大星相近. 为便于研究,如图,以大星的中心点为原点,建立直角坐标系, ${{{OO}}}_{1},{{{OO}}}_{2},{{{OO}}}_{3},{{{OO}}}_{4}$ 分别是大星中心点与四颗小星中心点的连接线，$a \approx  {16}^{ \circ  }$ ,求第三颗小星的一条边 ${AB}$ 所在直线的倾斜角?
    
    \begin{center}
    \includegraphics[width=0.3\textwidth]{bodyxb/src/bo_d2nqu4ref24c73b1hs1g_3_239_1048_362_323_0.jpg}
    \end{center}
    \begin{jieda}
        $\because O,{O}_{3}$ 都为五角星的中心点,
        
        $\therefore O{O}_{3}$ 平分第三颗小星的一个角,
        
        又五角星的内角为 ${36}^{ \circ  }$ 知: $\angle {BA}{O}_{3} = {18}^{ \circ  }$ ,
        
        过 ${O}_{3}$ 作 $x$ 轴的平行线 ${O}_{3}E$ ,如下图,则 $\angle O{O}_{3}E = \angle a \approx  {16}^{ \circ  }$ ,
        
        \begin{center}
        \includegraphics[width=0.3\textwidth]{bodyxb/src/bo_d2ntcqref24c73b1ilq0_1_242_1703_362_321_0.jpg}
        \end{center}
        
        $\therefore$ 直线 ${AB}$ 的倾斜角为 ${18}^{ \circ  } - {16}^{ \circ  } = {2}^{ \circ  }$ .
    \end{jieda}
\end{enumerate}


\clearpage
\section*{直线的方程}
\subsection*{直线的方程 (1)}

\zsd 直线的点斜式方程、斜截式方程、直线的两点式方程、点法式方程、一般式方程

补充:点方向式方程、截距式方程.

A 组：
\begin{enumerate}
    \item 经过点 $P\left( {-1,2}\right)$ 且斜率为 2 的直线 $l$ 的点斜式方程是\blkda[4]{$y - 2 = 2\left( {x - \left( {-1}\right) }\right)$}.
    
    \item 经过点 $P\left( {-1,2}\right)$ 且 ${y}$ 轴的截距为 -2 的直线 $l$ 的斜截式方程是\blkda[3]{$y = -4x-2$}.
    
    \item 经过点 ${A}\left( {-1,0}\right)$ , ${B}\left( {0,2}\right)$ 的直线 $l$ 的截距式方程是\blkda[3]{$\dfrac{x}{-1} + \dfrac{y}{2} = 1$}.

    \item 经过 $M\left( {1,2}\right) ,N\left( {5,3}\right)$ 的两点式方程是\blkda[3]{$\dfrac{y - 2}{3-2} = \dfrac{x - 1}{5-1}$}.

    \item 过点$(-1,0)$,倾斜角为 $\pi  - \arctan \dfrac{3}{4}$ 的直线的点方向式方程为\blkda[4]{$\dfrac{x - \left( {-1}\right) }{4} = \dfrac{y - 0}{-3}$}.
    \begin{jieda}
        若倾斜角为 $\pi  - \arctan \dfrac{3}{4}$ ,则直线的斜率 $k =  - \dfrac{3}{4}$ ,

        则直线的一个方向向量为 $\vv{d} = \left( {4, - 3}\right)$ ,又直线过点$(-1,0)$,

        故直线的点方向式方程为 $\dfrac{x - \left( {-1}\right) }{4} = \dfrac{y - 0}{-3}$ .

    \end{jieda}

    \item 写出下列直线的一个方向向量 $\vv{d}$ 和法向量 $\vv{n}$ :
    \begin{dingautolist}{172}
        \item 直线 $y + 2 = 0$ 的 $\vv{d} =$ \blkda{$(1, 0)$}, $\vv{n} =$ \blkda{$(0, 1)$}；
        \item 直线 $x - 3 = 0$ 的 $\vv{d} =$ \blkda{$(0, 1)$}, $\vv{n} =$ \blkda{$(1, 0)$}；
    \end{dingautolist}
    
    \item 若直线 ${4x} + {3y} - 5 = 0$ 的一个法向量为 $\left( {a + 2,a}\right)$ ,则实数 $a$ 的值为\blkda{6}.
    \begin{jieda}
        一个法向量$\vv{n} = (4, 3)$，故$\dfrac{a+2}{4} = \dfrac{a}{3}$，解得$a = 6$.
    \end{jieda}

\end{enumerate}
B 组：
\begin{enumerate}
    \item 直线经过点 $M\left( {-2, - 1}\right)$ 且在 $x$ 轴和 $y$ 轴的截距相等,则直线 $l$ 的斜截式方程是\blkda[4]{$y = -x-3$或$y = \dfrac{x}{2}$}.
    \begin{jieda}
        \begin{dingautolist}{172}
            \item 若截距为$0$，则设直线方程为$y = kx$，带入$M\left( {-2, - 1}\right)$可得$y = \dfrac{x}{2}$
            \item 若截距不为$0$，则设直线方程为$\dfrac{x}{a} + \dfrac{y}{a} = 1$，带入$M\left( {-2, - 1}\right)$可得$a = -3$，故$y = -x-3$
        \end{dingautolist}
    \end{jieda}

    \item 已知三点 $A\left( {4,2}\right) ,B\left( {-1, - 3}\right) ,P\left( {1,5}\right)$ ,若点 $Q$ 为线段 ${AB}$ 上的动点,则直线 ${PQ}$ 的斜率的范围为\blkda[4]{$(-\infty, -1] \cup [4, +\infty)$}.

    \item 若直线 $l$ 的倾斜角为 $\arccos \dfrac{4}{5}$ ,且它与坐标轴围成的三角形的面积为 6,则直线 $l$ 的截距式方程为\blkda[5]{$\dfrac{x}{4} + \dfrac{y}{-3} = 1$或$\dfrac{x}{-4}+\dfrac{y}{3} = 1$}.

    \item 已知直线 ${l}_{1}$ 过点 $P\left( {2,1}\right)$ ,且直线 ${l}_{1}$ 的倾斜角是直线 ${l}_{2} : {3x} - {4y} - 8 = 0$ 的倾斜角的一半,求直线 ${l}_{1}$ 的斜截式方程\blkda[3]{$y = \dfrac{1}{3}x+\dfrac{1}{3}$}.

    \item 设直线 ${2x} + \left( {k - 3}\right) y - {2k} + 6 = 0$ 过定点 $P$ ,则点 $P$ 的坐标为\blkda{$\left( {0,2}\right)$}.
    \begin{jieda}
        由直线方程 ${2x} + \left( {k - 3}\right) y - {2k} + 6 = 0$ ,可化简为 $\left( {{2x} - {3y} + 6}\right)  + k\left( {y - 2}\right)  = 0$ ,

        又由 $\left\{  \begin{array}{l} {2x} - {3y} + 6 = 0 \\  y - 2 = 0 \end{array}\right.$ ,解得 $x = 0,y = 2$，即直线恒经过定点 $P\left( {0,2}\right)$ .
    \end{jieda}

    \item 已知直线 ${l}_{1}$ 与直线 ${l}_{2} : {2x} - {3y} + 4 = 0$ 平行,且直线 ${l}_{1}$ 在 $y$ 轴上的截距为 -2,则直线 ${l}_{1}$ 的斜截式方程是\blkda[3]{$y = \dfrac{2}{3}x-2$}.
    
    \item 已知直线 ${3x} + {4y} = b$ 与两坐标轴围成的三角形的面积为 $\dfrac{3}{2}$ ,则 $b =$ \blkda{$6$ 或 $-6$}.
    \begin{jieda}
        令 $x = 0$ ,得 $y = \dfrac{b}{4}$ ; 令 $y = 0$ ,得 $x = \dfrac{b}{3}$ .
        
        故与坐标轴围成的三角形的面积为 $S = \dfrac{1}{2} \times  \left| \dfrac{b}{4}\right|  \times  \left| \dfrac{b}{3}\right|  = \dfrac{3}{2}$ ,解得 $b =  \pm  6$
    \end{jieda}
    
    \item 方程 $y - 1 = k\left( {x - 0}\right)$ 表示\dotfill{(\kh{C})}
    \xx{通过点$(0,1)$的一切直线}{通过点$(0, -1)$的一切直线}{过点$(0,1)$且不垂直于 $x$ 轴的一切直线}{过点$(0, -1)$且不垂直于 $x$ 轴的一切直线}
    
    \item 若点 $A\left( {a + b,{ab}}\right)$ 在第二象限内,则直线 ${bx} + {ay} - {ab} = 0$ 不经过\dotfill{(\kh{A})}
    \xx{第一象限}{第二象限}{第三象限}{第四象限}
    
    \item 已知直线 $l$ 被两条直线 ${l}_{1} : {2x} + y + 3 = 0$ 和 ${l}_{2} : {3x} - {5y} + 1 = 0$ 截得的线段的中点为 $P\left( {\dfrac{2}{3}, - 3}\right)$ ,则直线 $l$ 的一般式方程为\blkda[4]{${60x} + {17y} + {11} = 0$}.
    \begin{jieda}
        设 $l$ 与 ${l}_{1}$ 的交点为 $\left( {{x}_{0},{y}_{0}}\right)$ ,则 $l$ 与 ${l}_{2}$ 的交点为 $\left( {\dfrac{4}{3} - {x}_{0}, - 6 - {y}_{0}}\right)$
        
        所以有 $\left\{ \begin{array}{l}
             2{x}_{0} + {y}_{0} + 3 = 0  \\
             3\left( {\dfrac{4}{3} - {x}_{0}}\right)  - 5\left( {-6 - {y}_{0}}\right)  + 1 = 0
        \end{array}\right.$
        
        联立解得 $\left\{  \begin{array}{l} {x}_{0} = \dfrac{20}{13} \\  {y}_{0} =  - \dfrac{79}{13} \end{array}\right.$，所以 $l : \dfrac{y + \dfrac{79}{13}}{-3 + \dfrac{79}{13}} = \dfrac{x - \dfrac{20}{13}}{\dfrac{2}{3} - \dfrac{20}{13}}$ ,整理得 ${60x} + {17y} + {11} = 0$ .
    \end{jieda}
    
    \item 已知 $A\text{ 、 }B$ 是 $x$ 轴上的两点,点 $P$ 的横坐标为 2,直线 ${PA}$ 的方程为 $x - y + 1 = 0$ ,且 $\left| {PA}\right|  = \left| {PB}\right|$ ,求直线 ${PB}$ 的斜截式方程.
    \begin{jieda}
        $P(2, 3)$，$A(-1, 0)$，因为$\left| {PA}\right|  = \left| {PB}\right|$，所以$B(5, 0)$，故$PB: y = -x+5$
    \end{jieda}
    
    \item 已知直线 $l$ 经过点 $P\left( {-1,2}\right)$ .
    \begin{enumerate}
        \item 若 $l$ 不过原点且在两坐标轴上截距和为零,求 $l$ 的点斜式方程；
        \item 设 $l$ 的斜率 $k > 0$ , $l$ 与两坐标轴的交点分别为 ${A}$ 、 $B$ ,当 ${\Delta AOB}$ 的面积最小时，求 $l$ 的斜截式方程.
    \end{enumerate}
    \begin{jieda}
        \begin{enumerate}
            \item 由题意知, $l$ 的斜率存在且不为 0，设斜率为 $k$ ,则 $l$ 的点斜式方程为 $y - 2 = k\left( {x + 1}\right)$ ,则它在两坐标轴上截距分别为 $- 1 - \dfrac{2}{k}$ 和 $k + 2$，所以 $- 1 - \dfrac{2}{k} + k + 2 = 0$，解得 $k =  - 2$ (此时直线过原点,舍去) 或 $k = 1$，所以 $l$ 的点斜式方程为 $y - 2 = x + 1$ .
            \item 由(1)知, $A\left( {-\dfrac{2}{k} - 1,0}\right) ,B\left( {0,k + 2}\right) ,\left( {k > 0}\right)$，所以 $\triangle{AOB}$ 的面积 $S = \dfrac{1}{2}\left| {-\dfrac{2}{k} - 1}\right|  \cdot  \left| {k + 2}\right|  = \dfrac{{\left( k + 2\right) }^{2}}{2k} = \dfrac{k}{2} + 2 + \dfrac{2}{k} \geq  2\sqrt{\dfrac{k}{2} \cdot \dfrac{2}{k}} + 2 = 4$ ,

            当且仅当 $\dfrac{k}{2} = \dfrac{2}{k}$ 即 $k = 2$ 时,等号成立, $l$ 的点斜式方程为 $y - 2 = 2\left( {x + 1}\right)$ ,即 $y = {2x} + 4$
        \end{enumerate}
    \end{jieda}
    
    \item \begin{enumerate}
        \item 过点 $P\left( {2,1}\right)$ 作直线 $l$ 交 $x$ 轴, $y$ 轴的正方向于 $A,B$ ,求当 $\bigtriangleup {AOB}$ 面积最小时，直线 $l$ 的方程
        \item 已知点 $P\left( {a,b}\right)$ 为第一象限内一定点,过 $P$ 的动直线分别交 $x$ 轴、$y$ 轴的正半轴于点 $M\text{ 、 }N$ . 求证: 当 $P$ 为线段 ${MN}$ 的中点时, $\bigtriangleup {MON}$ 的面积为最小
    \end{enumerate}
    \begin{jieda}
        \begin{enumerate}
            \item 设$l: y = k(x-2)+1$（$k < 0$），$B(0, 1-2k), A\left(2-\dfrac{1}{k}, 0\right)$，故$S = \dfrac{1}{2}(1-2k)\left(2-\dfrac{1}{k}\right) = \dfrac{4-4k-\dfrac{1}{k}}{2} \geq 4$，当且仅当$k = -\dfrac{1}{2}$时取得最值.

            故直线方程为$x+2y-4=0$
            \item 设$l: y = k(x-a)+b$（$k < 0$），$N(0, b-ak), M\left(a-\dfrac{b}{k}, 0\right)$，故$S = \dfrac{1}{2}(b-ak)\left(a-\dfrac{b}{k}\right) = \dfrac{2ab-a^2k-\dfrac{b^2}{k}}{2} \geq 2ab$，当且仅当$k = -\dfrac{b}{a}$时取得最值，此时$M(2a, 0), N(0, 2b)$，$P$为$AB$中点.
        \end{enumerate}
    \end{jieda}
    \item 已知直线 $l : y = {3x}$ ,求过点 $A\left( {3,2}\right)$ 的直线 $m$ ,使它和直线 $l$ 、 $x$ 轴在第一象限内围成的三角形的面积最小.
    \begin{jieda}
        \begin{dingautolist}{172}
            \item 当斜率不存在时，直线方程为$x = 3$，此时三角形面积为$\dfrac{27}{2}$
    
            \item 当斜率存在时，设$y = k(x-3)+2$，由题意易知$k \in (-\infty, 0) \cup (3, +\infty)$
        
                设$l$与$m$交点为$B$，与$x$轴交点为$C$，则$S = \dfrac{1}{2} \times y_B \times x_C = \dfrac{1}{2} \times \dfrac{3(2-3k)}{3-k} \times \left(-\dfrac{2-3k}{k}\right) = \dfrac{3(2-3k)^2}{2(k-3)k}$
        
                设$b = 2-3k, b \in (-\infty, -7) \cup (2, +\infty)$，则$S = \dfrac{27b^2}{2(7+b)(b-2)} = \dfrac{27}{2} \times \dfrac{1}{1+\dfrac{5}{b}-\dfrac{14}{b^2}}$，根据二次函数性质可知其在$b = \dfrac{28}{5}$时取得最小值$\dfrac{28}{3}$，因此$k = -\dfrac{6}{5}$
        \end{dingautolist}
        

        综上，目标直线方程为$6x+5y-28=0$.
    \end{jieda}

\end{enumerate}

\clearpage
\subsection*{直线的方程 (2)}
\begin{enumerate}
    \item 过点 $P\left( {\cos a,\sin a}\right) \left( {-\dfrac{\pi }{2} < a < 0}\right)$ ,且以 $\vv{OP}$ ( $O$ 为平面直角坐标系的原点) 为法向量的直线的一般式方程为\blkda[5]{$\cos a \cdot x + \sin a \cdot y - 1 = 0$},它的倾斜角为\blkda{$\dfrac{\pi}{2} + a$}.
    \begin{jieda}
        根据法向量设直线方程为$\cos a \cdot x + \sin a \cdot y + c = 0$，代入$P\left( {\cos a,\sin a}\right)$坐标，根据平方关系可得$c = -1$，故直线的一般式方程为$\cos a \cdot x + \sin a \cdot y - 1 = 0$
        
        斜率为$-cota$，因为$-\dfrac{\pi }{2} < a < 0$，所以$-cota > 0$，根据诱导公式可知$-cota = tan\left(\dfrac{\pi}{2}+a\right)$倾斜角为$\dfrac{\pi}{2} + a$.
    \end{jieda}
    
    \item 已知原点在直线 $l$ 的射影为点$(-1,2)$,则直线 $l$ 的点法式方程是\blkda[4]{$-(x+1)+2(y-2) = 0$}.
    \begin{jieda}
        法向量$\vv{n} = (-1,2)$，直线点法式方程为$-(x+1)+2(y-2) = 0$
    \end{jieda}
    
    \item 若直线 $l$ 过点 $A\left( {5, - 4}\right)$ ,且它在两坐标轴上的截距互为相反数,则直线 $l$ 的方程为\blkda[4]{$y = -\dfrac{4}{5}x$或$y = x-9$}.

    \item 已知直线 ${l}_{1}$ 与直线 ${l}_{2} : {3x} - {4y} - 7 = 0$ 垂直,且直线 ${l}_{1}$ 与两坐标轴所构成的三角形的周长为 10 ,则线 ${l}_{1}$ 的方程是\blkda[4]{$4x+3y\pm10 = 0$}.
    \begin{jieda}
        根据垂直条件可设直线$l_1: 4x+3y + c = 0$与坐标轴交点为$\left(0, -\dfrac{c}{3}\right), \left(-\dfrac{c}{4}, 0\right)$，故$\left|\dfrac{c}{3}\right|+\left|\dfrac{c}{4}\right|+\left|\dfrac{5c}{12}\right| = 10$，因此$c = \pm 10$
    \end{jieda}

    \item 已知线段 ${PQ}$ 两端点的坐标分别为 $\left( {-1,1}\right) \text{ 、 }\left( {2,2}\right)$ ,若直线 $l : x + {my} + m = 0$ 与线段 ${PQ}$ 有交点,则 $m$ 的取值范围是\blkda{$\left[-\dfrac{2}{3}, \dfrac{1}{2}\right]$}.
    \begin{jieda}
        直线$l: \left\{\begin{array}{ll}
             x = 0, & m = 0 \\
             y = -\dfrac{x}{m}-1, & m \neq 0
        \end{array}\right.$，数形结合可得$m\in \left[-\dfrac{2}{3}, \dfrac{1}{2}\right]$
    \end{jieda}

    \item 已知 $4{x}_{1} - 3{y}_{1} + 5 = 0,4{x}_{2} - 3{y}_{2} + 5 = 0$ . 求过 $A\left( {{x}_{1},{y}_{1}}\right) \text{ 、 }B\left( {{x}_{2},{y}_{2}}\right)$ 两点的直线方程\blkda[3]{$4{x} - 3{y} + 5 = 0$}.

    \item 已知直线 ${l}_{1} : \left( {2 - a}\right) x + {ay} + 3 = 0$ 和直线 ${l}_{2} : x - {ay} - 3 = 0$ . “ $a = 1$ ” 是 “直线 ${l}_{1}$ 的法向量恰是直线 ${l}_{2}$ 的方向向量” 的\dotfill{(\kh{A})}
    \xx{充分非必要条件}{必要非充分条件}{充要条件}{既非充分又非必要条件}
    \begin{jieda}
        “直线 ${l}_{1}$ 的法向量恰是直线 ${l}_{2}$ 的方向向量”即两直线的法向量相互垂直，$\vv{n_1} = (2-a, a)$，$\vv{n_2} = (1, -a)$，因此有$2-a-a^2 = 0$，解得$a = 1$或$-2$.
    \end{jieda}
    
    \item 求过点 $A\left( {-1,1}\right)$ ,且与点 $B\left( {2,5}\right)$ 的距离最大的直线 $l$ 的点法向式方程\blkda[4]{$3(x+1)+4(y-1) = 0$}.
    \begin{jieda}
        根据题意直线应与向量$\vv{AB} = (3, 4)$垂直，故$\vv{AB}$为法向量，因此点法式方程为$3(x+1)+4(y-1) = 0$
    \end{jieda}
    
    \item 若 $\bigtriangleup {ABC}$ 的三个顶点的坐标分别是 $A\left( {3, - 1}\right) ,B\left( {-5,3}\right) ,C\left( {1, - 5}\right)$ ,则 $\angle A$ 的角平分线所在的直线方程是\blkda[3]{$x-3y-6 = 0$}.
    \begin{jieda}
        $\vv{AB} = (-8, 4)$，$\vv{AC} = (-2, -4)$，$|\vv{AB}| = 2|\vv{AC}|$，故$\vv{AB}+2\vv{AC} = (-12, -4)$为$\angle{A}$角分线的方向向量. 故直线方程为$\dfrac{x-3}{-12} = \dfrac{y+1}{-4}$，即$x-3y-6 = 0$.
    \end{jieda}

    \item 在平面直角坐标系中,设 $M\left( {{x}_{1},{y}_{1}}\right) ,N\left( {{x}_{2},{y}_{2}}\right)$ 为不同的两点,直线 $I$ 的方程为 ${ax} + {by} + c = 0$ ,设 $\delta  = \dfrac{a{x}_{1} + b{y}_{1} + c}{a{x}_{2} + b{y}_{2} + c}$ ,其中 $a,b,c$ 均为实数,下列四个说法中:
    \begin{dingautolist}{172}
        \item 存在实数 $\delta$ ,使点 $N$ 在直线 $l$ 上
        \item 若 $\delta  = 1$ ,则过 $M,N$ 两点的直线与直线 $l$ 重合
        \item 若 $\delta  =  - 1$ ,则直线 $l$ 经过线段 ${MN}$ 的中点
    \end{dingautolist}
    所有结论正确的说法的序号是\blkda{\ding{173}\ding{174}}.
    \begin{jieda}
        \begin{dingautolist}{172}
            \item 点 $N$ 在直线 $l$ 上$\Longleftrightarrow $分母为0，故不可能.
            \item $\delta  = 1 \Longleftrightarrow a(x_1-x_2)+b(y_1-y_2) = 0$，故过 $M,N$ 两点的直线与直线 $l$ 平行（法向量相同），但点 $N$ 不在直线 $l$ 上，故必然是平行关系.
            \item $\delta  = -1 \Longleftrightarrow a(x_1+x_2)+b(y_1+y_2) + 2c = 0 \Longleftrightarrow a \cdot \dfrac{x_1+x_2}{2} + b \cdot \dfrac{y_1+y_2}{2} + c = 0$，故直线 $l$ 经过线段 ${MN}$ 的中点.
        \end{dingautolist}
    \end{jieda}
    
    \item 已知 $a > 0$ ,直线 ${l}_{1} : x + {ay} = {2a} + 4$ 与 $y$ 轴的交点为 ${A},{l}_{2} : {2x} + {ay} = {2a} + 8$ 与 $x$ 轴的交点为 $B$ , ${l}_{1}$ 与 ${l}_{2}$ 的交点为 $C$ ,则四边形 ${OACB}$ 的面积的最小值为\blkda{$8+4\sqrt{2}$}.
    \begin{jieda}
        因为直线 ${l}_{1} : x - 4 =  - a\left( {y - 2}\right) ,{l}_{2} : 2\left( {x - 4}\right)  =  - a\left( {y - 2}\right)$ 都过点$(4,2)$，所以点 $C$ 的坐标是$(4,2)$.
        
        在 $x + {ay} = {2a} + 4$ 中,令 $x = 0$ ,得 $y = 2 + \dfrac{4}{a}$ ,所以 $A\left( {0,2 + \dfrac{4}{a}}\right)$ ,
        
        在 ${2x} + {ay} = {2a} + 8$ 中,令 $y = 0$ ,得 $x = 4 + a$ ,所以 $B\left( {4 + a,0}\right)$ ,
        
        所以 ${S}_{OACB} = {S}_{\bigtriangleup {OAC}} + {S}_{\bigtriangleup {OCB}} = \dfrac{1}{2}\left( {2 + \dfrac{4}{a}}\right)  \times  4 + \dfrac{1}{2}\left( {4 + a}\right)  \times  2 = 8 + \left( {a + \dfrac{8}{a}}\right)  \geq  8 + 2\sqrt{a} \cdot  \dfrac{8}{a} = 8 + 4\sqrt{a}$ ,当且仅当 $a = \dfrac{8}{a}$ ,即 $a = 2\sqrt{2}$ 时等号成立.
        
        所以当 $a = 2\sqrt{2}$ 时,四边形 ${OACB}$ 的面积取最小值为 $8 + 4\sqrt{2}$ .
        
        \begin{center}
        \includegraphics[width=0.3\textwidth]{bodyxb/src/bo_d2ntiebef24c73b1imj0_4_240_194_367_307_0.jpg}
        \end{center}
    \end{jieda}
    
    \item 过点$(-1, - 4)$的直线 $l$ 分别与 $x$ 轴、 $y$ 轴交于不同的 $A,B$ 两点, $O$ 为坐标原点, 若存在 4 条直线 $l$ 使得 $\Delta {{AOB}}$ 的面积均为 $m$ ,则 $m$ 的取值范围是\blkda{$(8, +\infty)$}.
    \begin{jieda}
        \begin{enumerate}
            \item 方法1 \\
            设直线$l$方程为：$\dfrac{x}{a} + \dfrac{y}{b} = 1$，代入$(-1, - 4)$,则 $\dfrac{1}{a} + \dfrac{4}{b} =  - 1$，即 $a =  - \dfrac{b}{b + 4}$
            
            $S = \dfrac{1}{2}|a||b| = \dfrac{1}{2}\left|\dfrac{{b}^{2}}{b + 4}\right| = m$，故 ${b}^{2} = {2m}\left| {b + 4}\right|$ 有4个解，将等式两侧看做函数图象可知，若两个函数图像有四个交点时，必然有两个交点在$x = -4$左侧，两个交点在$x = -4$右侧，因此原方程有四个解即 ${b}^{2} = {2mb} + {8m}$ 和 ${b}^{2} =  - {2mb} - {8m}$ 各2个解
            
            $\left\{ \begin{array}{l}
                 {\Delta }_1 = 4\left( {{m}^{2} + {8m}}\right)  > 0  \\
                 {\Delta }_2 = 4\left( {{m}^{2} - {8m}}\right)  > 0 
            \end{array}\right.$
            又 $m > 0$ ,故 $m \in  \left( {8, + \infty }\right)$

            \item 显然直线 ${AB}$ 的斜率存在且不为 0,设直线 ${AB}$ 的方程为 $y + 4 = k\left( {x + 1}\right)$ ,
            
            令 $x = 0$ 得 $y = k - 4$ ,令 $y = 0$ 得 $x = \dfrac{4}{k} - 1$ ,
            
            则 ${S}_{\bigtriangleup {OAB}} = \dfrac{1}{2}\left| {k - 4}\right| \left| {\dfrac{4}{k} - 1}\right|  = \dfrac{1}{2}\left| {8 - k - \dfrac{16}{k}}\right|$ ,
            
            由题意关于 $k$ 的方程 $\dfrac{1}{2}\left| {8 - k - \dfrac{16}{k}}\right|  = m$ 有四个不同的实数解，即$8 - k - \dfrac{16}{k} =  \pm  {2m}$

            结合对勾函数性质可知四根必定不重复，所以等价于${k}^{2} + \left( {{2m} - 8}\right) k + {16} = 0$ 有两个不等实根且 ${k}^{2} - \left( {{2m} + 8}\right) k + {16} = 0$ 有两个不等实根,
            
            $\left\{  \begin{array}{l} {\Delta }_{1} = {\left( 2m - 8\right) }^{2} - {64} > 0 \\  {\Delta }_{2} = {\left( 2m + 8\right) }^{2} - {64} > 0 \end{array}\right.$ ,解得 $m <  - 8$ 或 $m > 8$ .
            
            又 $m > 0$ ,所以 $m > 8$ .
        \end{enumerate}
        
    \end{jieda}
    
    \item 已知直线 $l$ 过点 $P\left( {1,a}\right)$ ,且直线 $l$ 与 $x$ 轴、 $y$ 轴分别交于 $A\text{ 、 }B$ 两点,若 $P$ 恰为 ${AB}$ 的中点,求直线 $l$ 的斜率和倾斜角.
    \begin{jieda}
        $y = -ax+2a$，斜率为$-a$
        \begin{dingautolist}{172}
            \item $a > 0$，倾斜角为$\pi - arctan(a)$
            \item $a < 0$，倾斜角为$-arctan(a)$
        \end{dingautolist}
    \end{jieda}
    
    \item 过定点 $M\left( {0,1}\right)$ 作一直线 $l$ ,使它夹在两已知直线 ${l}_{1} : x - {3y} + {10} = 0$ 和 ${l}_{2} : {2x} + y - 8 = 0$ 之间的线段被点 $M$ 平分,求直线 $l$ 的方程.
    \begin{jieda}
        设$l$与$l_2, l_1$交于$A,B$
    
        根据点 $A$ 在直线 ${l}_{2}$ 上,可设 $A\left( {a,8 - {2a}}\right)$，因为$M\left( {0,1}\right)$ 是 ${AB}$ 的中点，所以$B\left( {-a,{2a} - 6}\right)$ 
        
        由于点 $B$ 在直线 ${l}_{1}$ 上，故有$- a - 3\left( {{2a} - 6}\right)  + {10} = 0$，解得: $a = 4$
        
        即 $A\left( {4,0}\right)$ ,直线 $l$ 的斜率为 ${k}_{MA} =  - \dfrac{1}{4}$，则直线 $l$ 的方程为 $y - 0 =  - \dfrac{1}{4}\left( {x - 4}\right)$ ,即 $x + {4y} - 4 = 0$.
    \end{jieda}
    
    \item 如图,在直角梯形 ${ABCD}$ 中, $\angle C = \angle D = {90}^{ \circ  }$ ,两平行边 $\dot{A}D$ 、 ${BC}$ 上各有一个动点 $P$ 、 $Q$ ,直线 ${PQ}$ 把梯形 ${ABCD}$ 分成面积相等的两部分,直线 ${PQ}$ 是否过定点? 若过定点, 试建立适当的平面直角坐标系, 并求出定点的坐标; 若不过定点, 说明理由.
    
    \begin{center}
    \includegraphics[width=0.2\textwidth]{bodyxb/src/bo_d2nqucf7aajc738rmrjg_2_236_497_304_267_0.jpg}
    \end{center}
    \begin{jieda}
        以$D$为原点建系，设$A(a, 0), C(0, c), B(b, c)$

        设$P(p, 0), Q(q, c)$，直线为$y = \dfrac{c}{q-p}(x-p)$
        
        根据面积关系有$\dfrac{(p+q)c}{2} \times 2 = \dfrac{(a+b)c}{2}$，因此$p+q = \dfrac{a+b}{2}$，代入直线方程有$y = \dfrac{c}{\dfrac{a+b}{2}-2p}(x-p)$，过定点$\left(\dfrac{a+b}{4}, \dfrac{c}{2}\right)$
    \end{jieda}

\end{enumerate}

\clearpage
\section*{两条直线的位置关系}
\subsection*{两条直线的位置关系 (1)}

\zsd 两直线位置关系判定的方法、两直线垂直、平行、重合的充要条件

A 组：
\begin{enumerate}
    \item 已知直线 ${l}_{1} : {mx} + y - m - 1 = 0$ 和直线 ${l}_{2} : x + {my} - {2m} = 0$ . 当 $m =$ \blkda{$-1$}时, ${l}_{1}//{l}_{2}$ ; 当 $m =$ \blkda{$1$}时, ${l}_{1}$ 与 ${l}_{2}$ 重合.

    \item 直线 ${l}_{1} : {px} + {3y} + 1 = 0$ 与直线 ${l}_{2} : {6x} - {2y} - 5 = 0$ 垂直,则 $p$ 的值为\blkda{$1$}.
    \begin{jieda}
        $6p-3\times(-2) = 0$，解得$p = 1$
    \end{jieda}

    \item 已知直线 ${l}_{1} : x + {my} + 6 = 0,{l}_{2} : \left( {m - 2}\right) x + {3y} + {2m} = 0$ ,求实数 $m$ 的值,使得直线 ${l}_{1}$ 和 ${l}_{2}$ :
    \begin{enumerate*}
        \item 相交
        \item 垂直
        \item 平行
        \item 重合
    \end{enumerate*}
    \begin{jieda}
        \begin{enumerate*}
            \item $m \notin \{-1, 3\}$
            \item $m = \dfrac{1}{2}$
            \item $m = -1$
            \item $m = 3$
        \end{enumerate*}
    \end{jieda}

\end{enumerate}
B 组：
\begin{enumerate}
    \item 已知直线 ${l}_{1} : {ax} + \left( {a - 1}\right) y + 3 = 0,{l}_{2} : {2x} + {ay} - 1 = 0$ ,若 ${l}_{1} \bot  {l}_{2}$ ,则实数 $a$ 的值是\blkda{$-1$或$0$}.
    \begin{jieda}
        $2a+a(a-1) = 0$，即$a^2+a = 0$，解得$a = 0$或$-1$.
    \end{jieda}
    
    \item 若直线 ${2x} - y - {10} = 0$ 经过直线 ${4x} + {3y} - {10} = 0$ 和 ${ax} + {2y} + 8 = 0$ 的交点,则实数 $a =$ \blkda{$-1$}.
    \begin{jieda}
        直线 ${2x} - y - {10} = 0$ 和直线 ${4x} + {3y} - {10} = 0$交于$(4, -2)$，因此直线${ax} + {2y} + 8 = 0$过该点，故$a = -1$
    \end{jieda}
    
    \item 若直线 ${l}_{1} : y = {kx} + k + 2$ 与直线 ${l}_{2} : y =  - {2x} + 4$ 的交点在第一象限内,则实数 $k$ 的取值范围是\blkda{$\left(-\dfrac{2}{3}, 2\right)$}.
    \begin{jieda}
        ${kx} + k + 2 =  - {2x} + 4$，解得$x = \dfrac{2-k}{2+k}$，若交点在第一象限内，则有$\dfrac{2-k}{2+k} \in (0, 2)$，解得$k \in \left(-\dfrac{2}{3}, 2\right)$
    \end{jieda}
    
    \item 若三条直线 ${l}_{1} : x - y = 0,{l}_{2} : x + y - 2 = 0$ 和 ${l}_{3} : {5x} - {ky} - {15} = 0$ 能围成一个三角形, 则实数 $k$ 的取值范围是\blkda[4]{$\{k | k \neq \pm 5, k \neq -10\}$}.
    \begin{jieda}
        直线$l_1, l_2$交点为$(1, 1)$，当直线$l_3$过该点或与另两条直线平行时不能围成三角形，否则均可以.
    \end{jieda}
    
    \item 直线 ${l}_{1}$ 过点 $A\left( {m,1}\right)$ 和点 $B\left( {-1,m}\right)$ ,直线 ${l}_{2}$ 过点 $C\left( {m + n,n + 1}\right)$ 和点 $D\left( {n + 1,n - m}\right)$ ,则直线 ${l}_{1}$ 与 ${l}_{2}$ 的位置关系是\blkda{相交}.

    \item 直线 ${l}_{1} : x + \left( {m + 1}\right) y - {2m} - 2 = 0$ 与直线 ${l}_{2} : \left( {m + 1}\right) x - y - {2m} - 2 = 0$ 相交于点 $P$ ,对任意实数 $m$ ,直线 ${l}_{1},{l}_{2}$ 分别恒过定点 $A,B$ ,则 $\left| {PA}\right|  + \left| {PB}\right|$ 的最大值为\blkda{$4$}.
    \begin{jieda}
        因为 ${l}_{1} : x + \left( {m + 1}\right) y - {2m} - 2 = 0$ ,即 $x + y - 2 + m\left( {y - 2}\right)  = 0$，由 $\left\{  \begin{array}{l} x + y - 2 = 0 \\  y - 2 = 0 \end{array}\right.$ ,解得 $\left\{  \begin{array}{l} x = 0 \\  y = 2 \end{array}\right.$ ,所以直线 ${l}_{1}$ 过定点 $A\left( {0,2}\right)$ ;
        同理可得直线 ${l}_{2}$ 过定点 $B\left( {2,0}\right)$ ；
        
        又因为 $1 \times  \left( {m + 1}\right)  + \left( {-1}\right)  \times  \left( {m + 1}\right)  = 0$ ,所以 ${l}_{1} \bot  {l}_{2}$ ,
        即有 ${PA}\bot {PB}$ ,所以 ${\left| PA\right| }^{2} + {\left| PB\right| }^{2} = {\left| AB\right| }^{2} = 8$
        
        所以 $\left| {PA}\right|  + \left| {PB}\right|  \leq  \sqrt{2\left( {{\left| PA\right| }^{2} + {\left| PB\right| }^{2}}\right) } = 4$ ,
        当且仅当 $\left| {PA}\right|  = \left| {PB}\right|  = 2$ 时,取等号.
        所以 $\left| {PA}\right|  + \left| {PB}\right|$ 的最大值为 4 .
    \end{jieda}
    
    \item “ $a = 3$ ” 是 “直线 ${l}_{1} : \left( {a - 1}\right) x + {2y} + 1 = 0$ 与直线 ${l}_{2} : {3x} + {ay} - 1 = 0$ 平行” 的\dotfill{(\kh{C})}
    \xx{充分不必要条件}{必要不充分条件}{充要条件}{既不充分也不必要条件}
    
    \item “直线 ${l}_{1}$ 与直线 ${l}_{2}$ 平行” 是 “直线 ${l}_{1}$ 与直线 ${l}_{2}$ 的斜率相等” 的\dotfill{(\kh{D})}
    \xx{充分非必要条件}{必要非充要条件}{充要条件}{既非充分又非必要条件}
    \begin{jieda}
        非必要性：重合；非充分性：斜率不存在
    \end{jieda}
    
    \item 直线 ${l}_{1} : {3x} - {2y} + m = 0$ 和 ${l}_{2} : \left( {{m}^{2} + 1}\right) x + {3y} - {3m} = 0$ 的位置关系是\dotfill{(\kh{C})}
    \xx{平行}{重合}{相交}{不确定}
    \begin{jieda}
        $k_1 = \dfrac{3}{2}$，$k_2 = -\dfrac{m^2+1}{3}$，$k_1 > 0 > k_2$
    \end{jieda}
    
    \item 若 $a,b,c$ 是 $\bigtriangleup  {ABC}$ 的三边长,则直线 ${l}_{1} : x\sin A + {ay} + c = 0$ 与直线 ${l}_{2} : {bx} - y\sin B + \sin C = 0$ 的位置关系是\dotfill{(\kh{C})}
    \xx{平行}{重合}{垂直}{相交但不垂直}
    \begin{jieda}
        $k_1 = -\dfrac{sinA}{a} = -\dfrac{1}{2R}$，$k_2 = \dfrac{b}{sinB} = 2R$，故$k_1k_2 = -1$.
    \end{jieda}
    
    \item 设 $P\left( {{x}_{1},{y}_{1}}\right)$ 是直线 $l : f\left( {x,y}\right)  = 0$ 上一点, $Q\left( {{x}_{2},{y}_{2}}\right)$ 是 $l$ 外一点,则方程 $f\left( {x,y}\right)  = f\left( {{x}_{1},{y}_{1}}\right)  + f\left( {{x}_{2},{y}_{2}}\right)$ 表示的直线\dotfill{(\kh{C})}
    \xx{与 $l$ 重合}{与 $l$ 相交于 $P$ 点}{过 $Q$ 点且与 $l$ 平行}{过 $Q$ 点且与 $l$ 相交}
    
    \item 已知 $a > 0,b > 0,c > 0$ ,直线 ${l}_{1} : x\lg \left( {ac}\right)  - y + m = 0,{l}_{2} : x\lg \left( {bc}\right)  - y + n = 0$ 互相垂直, 求 $\dfrac{a}{b}$ 的取值范围.
    \begin{jieda}
        $l_1 \bot l_2 \Longleftrightarrow lg(ac) \times lg(bc) + 1 = 0  \Longleftrightarrow (lgb + lgc) (lga + lgc) + 1 = 0$，看做是关于$lgc$的二次方程$(lgc)^2 + (lga + lgb) lgc + (1+lgalgb) = 0$，要使得方程有解，则应有$\Delta = (lga + lgb)^2 - 4(1+lgalgb) \geq 0$，即$(lga - lgb)^2 \geq 4$，故$lg\left(\dfrac{a}{b}\right) \in (-\infty, -2] \cup [2, +\infty)$，即$a \in \left(0, \dfrac{1}{100}\right] \cup [100, +\infty)$
    \end{jieda}
    
    \item 已知直线 ${l}_{1} : \left( {m + 2}\right) x + \left( {{m}^{2} - {3m}}\right) y + 4 = 0$ 和直线${l}_{2} : {2mx} + 2\left( {m - 3}\right) y + m + 2 = 0$.
    \begin{enumerate}
        \item 当 $m$ 为何值时,直线 ${l}_{1}$ 与直线 ${l}_{2}$ 平行?
        \item 当 $m$ 为何值时,直线 ${l}_{1}$ 与直线 ${l}_{2}$ 重合?
    \end{enumerate}
    \begin{jieda}
        \begin{enumerate}
            \item $m \in \{-1, 3\}$
            \item $m = 2$
        \end{enumerate}
    \end{jieda}
    
    \item 已知集合 $A = \left\{  {\left( {x,y}\right) \left| {\dfrac{y - 3}{x - 2} = a + 1}\right. }\right\}$ ,集合 $B = \left\{  {\left( {x,y}\right) \left| {\left( {{a}^{2} - 1}\right) x + \left( {a - 1}\right) y - {15} = 0}\right. }\right\}$ . 是否存在实数 $a$ ,使 $A \cap  B = \varnothing$ ? 若存在,求出 $a$ 的值; 若不存在,说明理由.
    \begin{jieda}
        从形式上看$A, B$两个集合都是点集。\\
        集合$A$对应于直线$y = (a+1)(x-2) + 3$，但不包括其上的点$(2, 3)$ \\
        集合$B$对应于直线$y = -(a+1)x + \dfrac{15}{a-1} (a \neq 1)$，而$B = \varnothing (a = 1)$ 
        \begin{enumerate}
            \item 当$a = 1$时，由于$B = \varnothing$，自然满足条件
            \item 当$a \neq 1$时，若两直线平行或直线$y = -(a+1)x + \dfrac{15}{a-1}$过$(2, 3)$则满足条件
            \begin{dingautolist}{172}
                \item 当$a = -1$时，两直线平行
                \item 当$a = -4 \mbox{或} \dfrac{5}{2}$的时候直线$y = -(a+1)x + \dfrac{15}{a-1}$过$(2, 3)$
            \end{dingautolist}
        \end{enumerate}
        综上$a \in \left\{1, -1, -4, \dfrac{5}{2}\right\}$
    \end{jieda} 
\end{enumerate}

\clearpage
\subsection*{两条直线的位置关系 (2)}

\zsd 两直线的夹角公式

A组：
\begin{enumerate}
    \item 直线 ${l}_{1} : x + 2 = 0$ 与直线 ${l}_{2} : \sqrt{3}x + y + 3 = 0$ 的夹角的大小是\blkda{$\dfrac{\pi}{6}$}.
    
    \item 直线 ${l}_{1} : {2x} - {3y} + 1 = 0$ 与直线 ${l}_{2} : x - y - 3 = 0$ 的夹角的大小为\blkda{$arctan\dfrac{1}{5}$}.
    \begin{jieda}
        $k_1 = \dfrac{2}{3}, k_2 = 1$，$tan\theta = \dfrac{k_2-k_1}{1+k_1k_2} = \dfrac{1}{5}$
    \end{jieda}

    \item “两条直线的斜率的乘积等于 $-1$ ” 是 “这两条直线互相垂直” 的\blkda[3]{充分非必要}条件.
    
    \item 若直线 ${l}_{1} : \sqrt{3}x + y = 0$ 与直线 ${l}_{2} : {kx} - y + 1 = 0$ 的夹角的大小为 ${60}^{ \circ  }$ ,则实数 $k$ 的值为\blkda{$0$或$\sqrt{3}$}.

    \item 过点 $A\left( {1,\dfrac{\sqrt{3}}{3}}\right)$ ,且与直线 $x - \sqrt{3}y = 0$ 的夹角是 $\dfrac{\pi }{3}$ 的直线的方程是\blkda[4.5]{$x = 1$或$x+\sqrt{3}y-2 = 0$}.

    \item 已知直线 $l : \left( {m - 1}\right) x + y - {2m} - 9 = 0\left( {m \in  {R}}\right)$ ,则直线 $l$ 恒过定点\blkda{$(2, 11)$}.
    \begin{jieda}
        由题意可得 $m\left( {x - 2}\right)  - x + y - 9 = 0$ ,令 $\left\{  \begin{array}{l} x - 2 = 0 \\   - x + y - 9 = 0 \end{array}\right.$ ,解得 $x = 2,y = {11}$，所以直线 $l$ 恒过定点(2,11)
    \end{jieda}
\end{enumerate}

B 组：
\begin{enumerate}
    \item 直线 $l$ 过点$(1,2)$,且纵截距为横截距的两倍,则直线 $l$ 的方程是\blkda[5]{${2x} - y = 0$ 或 ${2x} + y - 4 = 0$}.
    \begin{jieda}
        \begin{dingautolist}{172}
            \item 当直线 $l$ 过原点时,直线斜率为 $\dfrac{2 - 0}{1 - 0} = 2$ ,直线方程为 $y = {2x}$ ,整理得 ${2x} - y = 0$ .
            \item 当直线 $l$ 不过原点时,设直线方程为 $\dfrac{x}{a} + \dfrac{y}{2a} = 1$ ,代入$(1,2)$得, $\dfrac{2}{a} = 1$ ,解得 $a = 2$ ,直线方程为 $\dfrac{x}{2} + \dfrac{y}{4} = 1$ ,整理得 ${2x} + y - 4 = 0$ .
        \end{dingautolist}
        综上，直线 $l$ 的方程为 ${2x} - y = 0$ 或 ${2x} + y - 4 = 0$ .
    \end{jieda}
    
    \item 已知等腰直角 $\bigtriangleup  {ABC}$ 的斜边所在直线的方程是 ${3x} - y + 2 = 0$ ,直角顶点是 $C\left( {3, - 2}\right)$ ,那么两条直角边 ${BC},{AC}$ 的方程分别为\blkda[6]{$x - {2y} - 7 = 0$ 或 ${2x} + y - 4 = 0$}.

     \item 已知直线 ${l}^{\prime }$ 经过点 $P\left( {2,1}\right)$ ,与直线 $l : x + {2y} + 1 = 0$ 的夹角为 $\arccos \dfrac{\sqrt{5}}{5}$ . 则直线 ${l}^{\prime }$ 的方程\blkda[4]{${3x} - {4y} - 2 = 0$或$x = 2$}.
    \begin{jieda}
        直线$l$的方向向量为$\vv{d} = (-2, 1)$，设$l'$的方向向量为$(cos\alpha, sin\alpha)$，则$cos\theta = \left|\dfrac{-2cos\alpha+sin\alpha}{\sqrt{5}}\right|$，根据题意可知$\left|\dfrac{-2cos\alpha+sin\alpha}{\sqrt{5}}\right| = \dfrac{\sqrt{5}}{5}$，故$\left|-2cos\alpha+sin\alpha\right| = 1$，即$3cos^2\alpha = 4sin\alpha cos\alpha$，故$tan\alpha = \dfrac{3}{4}$或$cos\alpha = 0$，即直线 ${l}^{\prime }$的斜率$k = \dfrac{3}{4}$或不存在.

        故直线${l}^{\prime }$ 的方程为${3x} - {4y} - 2 = 0$或 $x = 2$.
    \end{jieda}

    \item 直线 $y = {ax} - 2$ 与直线 $y = \sqrt{3}x$ 的夹角 $\theta  \in  \left\lbrack  {0,\dfrac{\pi }{6}}\right)$ ,则 $a$ 的取值范围是\blkda{$a > \dfrac{\sqrt{3}}{3}$}.
    \begin{jieda}
        由题知直线 $y = {ax} - 2$ 的斜率为 ${k}_{1} = a$ ,直线 $y = \sqrt{3}x$ 的斜率为 ${k}_{2} = \sqrt{3}$
        
        因为直线 $y = {ax} - 2$ 与直线 $y = \sqrt{3}x$ 的夹角 $\theta  \in  \left\lbrack  {0,\dfrac{\pi }{6}}\right)$，所以 $\tan \theta  = \left| \dfrac{{k}_{1} - {k}_{2}}{1 + {k}_{1}{k}_{2}}\right|  = \left| \dfrac{a - \sqrt{3}}{1 + \sqrt{3}a}\right|  \in  \left\lbrack  {0,\dfrac{\sqrt{3}}{3}}\right)$ ,即 $3{\left( a - \sqrt{3}\right) }^{2} < {\left( 1 + \sqrt{3}a\right) }^{2}$，解得 $a > \dfrac{\sqrt{3}}{3}$ .
    \end{jieda}

    \item 若入射光线所在的直线为 ${l}_{1} : x - y + 1 = 0$ 射入,被镜面直线 ${l}_{2} : {2x} + y - 4 = 0$ 反射,则反射光线所在的直线 $l$ 的方程是\blkda[3]{$x-7y+13=0$}.
    \begin{jieda}
        $l_1, l_2$的交点为$P(1,2)$
        
        在$l_1$上取点$M(0, 1)$，设$M$关于$l_2$的对称点为$M'(a, b)$，则$\vv{MM'} = (a, b - 1)$和$l_2$的方向向量$\vv{d} = (-1, 2)$垂直，且$MM'$的中点$N\left(\dfrac{a}{2}, \dfrac{b+1}{2}\right)$在直线$l_2$上，故$\left\{\begin{array}{l}
             -a+2(b-1) = 0  \\
             a + \dfrac{b+1}{2} - 4 = 0
        \end{array}\right.$，故$M'\left(\dfrac{12}{5}, \dfrac{11}{5}\right)$

        直线$l$即直线$M'P$，因此直线$l: x-7y+13=0$
    \end{jieda}

    \item 点 $P\left( {1, - 2}\right)$ 到直线 $l : \left( {3 + {2\lambda }}\right) x + \left( {4 + \lambda }\right) y - 2 + {2\lambda } = 0\left( {\lambda  \in  {R}}\right)$ 的距离最大时,其最大值以及此时的直线的方程分别为\blkda[4]{$5; {3x} - {4y} + {14} = 0$}.
    \begin{jieda}
        将直线 $l : \left( {3 + {2\lambda }}\right) x + \left( {4 + \lambda }\right) y - 2 + {2\lambda } = 0\left( {\lambda  \in  {R}}\right)$ 变形得 ${3x} + {4y} - 2 + \lambda \left( {{2x} + y + 2}\right)  = 0$
        
        由 $\left\{  \begin{array}{l} {3x} + {4y} - 2 = 0 \\  {2x} + y + 2 = 0 \end{array}\right.$ ,解得 $\left\{  \begin{array}{l} x =  - 2 \\  y = 2 \end{array}\right.$ ,因此直线 $l$ 过定点 $A\left( {-2,2}\right)$
        
        当 ${AP} \bot  l$ 时,点 $P\left( {1, - 2}\right)$ 到直线 $l : \left( {3 + {2\lambda }}\right) x + \left( {4 + \lambda }\right) y - 2 + {2\lambda } = 0\left( {\lambda  \in  \mathbf{R}}\right)$ 的距离最大，最大值为 $\left| {AP}\right|  = \sqrt{{\left( -2 - 1\right) }^{2} + {\left( 2 + 2\right) }^{2}} = 5$ ,又直线 ${AP}$ 的斜率 ${k}_{AP} = \dfrac{2 + 2}{-2 - 1} =  - \dfrac{4}{3}$，所以直线 $l$ 的方程为 $y - 2 = \dfrac{3}{4}\left( {x + 2}\right)$ ,即 ${3x} - {4y} + {14} = 0$ .
    \end{jieda}

    \item 已知直线 ${l}_{1}$ 过点 $M$ ,将直线 ${l}_{1}$ 绕点 $M$ 沿逆时针方向旋转 $a\left( {{0}^{ \circ  } < a < {90}^{ \circ  }}\right)$ 角,得到 的直线 ${l}_{2}$ 的方程是 $x + y - 2 = 0$ ,将直线 ${l}_{2}$ 绕点 $M$ 沿顺时针方向继续旋转 $\left( {{90}^{ \circ  } - a}\right)$ 角,得到的直线 ${l}_{3}$ 的方程是 ${2x} - y - 1 = 0$ ,则直线 ${l}_{1}$ 的方程是\blkda[3]{$2x+y-3 = 0$}.
    \begin{jieda}
        $l_2$和$l_3$的交点为$M$，故$M(1, 1)$

        $k_2 = -1, k_3 = 2, tan(\ang{90} - a) = \left|\dfrac{k_3-k_2}{1+k_2k_3}\right| = 3$，故$tan(a) = \dfrac{1}{3}$
        
        $\left|\dfrac{k_1-k_2}{1+k_1k_2}\right| = \dfrac{1}{3}$，解得$k_1 = -2$，因此直线 ${l}_{1}$ 的方程是$2x+y-3 = 0$
    \end{jieda}

    \item 已知点 $P$ 在直线 $y = 2$ 上运动,又已知有点 $A\left( {-2,0}\right) ,B\left( {2,0}\right)$ ,当 $\angle {APB}$ 最大时,则 $P$ 点坐标为\blkda{$(0, 2)$}.

    \item 已知直线 ${l}_{1} : y = x$ ,直线 ${l}_{2} : {ax} - y = 0\left( {a \in  \mathbf{R}}\right)$ . 当直线 ${l}_{1}$ 与直线 ${l}_{2}$ 的夹角在 $\left( {0,\dfrac{\pi }{12}}\right)$ 内变动时,则实数 $a$ 的取值范围是\blkda[3]{$\left(\dfrac{\sqrt{3}}{3}, 1\right) \cup (1, \sqrt{3})$}.

    \item 已知直线 ${l}_{1} : {mx} - y = 0\left( {m \in  {R}}\right)$ 过定点 ${A}$ ,直线 ${l}_{2} : x + {my} + 4 - {2m} = 0$ 过定点 $B,{l}_{1}$ 与 ${l}_{2}$ 的交点为 $C$ ,则 $\bigtriangleup  {{ABC}}$ 面积的最大值为\blkda{$5$}.
    \begin{jieda}
        由题可知, $A\left( {0,0}\right) ,B\left( {-4,2}\right)$ ,直线 ${l}_{1} \bot  {l}_{2}$，所以 ${AC}\bot {BC}$ , ${\left| AB\right| }^{2} = {20}$，所以 ${\left| AC\right| }^{2} + {\left| BC\right| }^{2} = {\left| AB\right| }^{2} = {20}$ ，所以 $\triangle {ABC}$ 的面积为 $\dfrac{1}{2}\left| {AC}\right| \left| {BC}\right|  \leq  \dfrac{1}{2} \times  \dfrac{{\left| AC\right| }^{2} + {\left| BC\right| }^{2}}{2} = 5$，当且仅当 $\left| {AC}\right|  = \left| {BC}\right|  = \sqrt{10}$ 时等号成立.
    \end{jieda}

    \item 已知 $a \in  \left( {0,2}\right)$ ,直线 ${l}_{1}$ 的方程为 ${ax} - {2y} - {2a} + 4 = 0$ ,直线 ${l}_{2}$ 的方程为 ${2x} + {a}^{2}y - 2{a}^{2} - 4 = 0$ ,记 ${l}_{1},{l}_{2}$ 与两坐标轴围成的四边形的面积为 $S$ ,则 $S$ 的最小值为\blkda{$\dfrac{15}{4}$}.
    \begin{jieda}
        由 ${ax} - {2y} - {2a} + 4 = 0 \Longrightarrow  2\left( {y - 2}\right)  = a\left( {x - 2}\right)  \Longrightarrow  y - 2 = \dfrac{a}{2}\left( {x - 2}\right)$，因此直线 ${l}_{1}$ 恒过点$(2,2)$,且该直线的斜率 ${k}_{1} \in  \left( {0,1}\right)$
        
        由 ${2x} + {a}^{2}y - 2{a}^{2} - 4 = 0 \Longrightarrow  {a}^{2}\left( {y - 2}\right)  =  - 2\left( {x - 2}\right)  \Longrightarrow  y - 2 =  - \dfrac{2}{{a}^{2}}\left( {x - 2}\right)$，因此直线 ${l}_{2}$ 也恒过点$(2,2)$,且该直线的 ${k}_{2} \in  \left( {-\infty , - \dfrac{1}{2}}\right)$

        设直线 ${l}_{1}$ 与纵轴的交点为 $A\left( {0,2 - a}\right)$ ,直线 ${l}_{2}$ 与横轴的交点为 $B\left( {2 + {a}^{2},0}\right)$，点$(2,2)$为 $C$ 点，所以 ${l}_{1}$ , ${l}_{2}$ 与两坐标轴围成的四边形为四边形 ${OACB}$
        如图所示:

        \begin{center}
            \includegraphics[width=0.3\linewidth]{bodyxb//src/zhixiandeweizhiguanxi1.png}
        \end{center}
        ${S}_{OACB} = {S}_{\bigtriangleup {AOC}} + {S}_{\bigtriangleup {OCB}} = \dfrac{1}{2}\left( {2 - a}\right)  \cdot  2 + \dfrac{1}{2} \cdot  \left( {2 + {a}^{2}}\right)  \cdot  2 = {a}^{2} - a + 4 = {\left( a - \dfrac{1}{2}\right) }^{2} + \dfrac{15}{4},$
        当 $a = \frac{1}{2}$ 时, ${S}_{OACB}$ 有最小值 $\dfrac{15}{4}$
    \end{jieda}

    \item 已知三条直线 ${l}_{1} : x - {2y} + 2 = 0,{l}_{2} : x - 2 = 0,{l}_{3} : x + {ky} = 0$ 将平面分为六个部分,则满足条件的 $k$ 的值共有\blkda{3}个.
    \begin{jieda}
        \begin{dingautolist}{172}
            \item 当三条直线交于一点时, 可将平面分为六个部分，联立 ${l}_{1} : x - {2y} + 2 = 0$ 与 ${l}_{2} : x - 2 = 0$ ,解得 $\left\{  \begin{array}{l} x = 2 \\  y = 2 \end{array}\right.$，则将 $\left\{  \begin{array}{l} x = 2 \\  y = 2 \end{array}\right.$ 代入 ${l}_{3} : x + {ky} = 0$ 中, ${2k} + 2 = 0$ ,解得 $k =  - 1$
            \item 当 ${l}_{3} : x + {ky} = 0$ 与 ${l}_{1} : x - {2y} + 2 = 0$ 平行时,满足要求,此时 $k =  - 2$
            \item 当 ${l}_{3} : x + {ky} = 0$ 与 ${l}_{2} : x - 2 = 0$ 平行时,满足要求,此时 $k = 0$
        \end{dingautolist}
        
        综上,满足条件的 $k$ 的值共有 3 个.
    \end{jieda}

    \item 已知直线 $m : y = x\cos a$ 和 $n : {3x} + y = c$ ,则\dotfill{(\kh{C})}
    \xx{$m$ 和 $n$ 可能重合}{$m$ 和 $n$ 不可能垂直}{存在直线 $m$ 上一点 $P$ ,以 $P$ 为中心旋转后与 $n$ 重合}{以上都不对}
    \begin{jieda}
        对 ${A}$ ,直线 $m : y = x\cos a$ ,斜率为 ${k}_{1} = \cos a$；直线 $n : {3x} + y = c$ ,斜率为 ${k}_{2} =  - 3$ ；
        ${k}_{1} \neq  {k}_{2}$ ,所以 $m$ 和 $n$ 不可能重合, ${A}$ 错误；
        
        对 ${B},\cos a = \dfrac{1}{3}$ 时, ${k}_{1} \times  {k}_{2} =  - 1,m$ 和 $n$ 垂直,所以 ${B}$ 错误;
        
        对 ${C}$ ,由 ${k}_{1} \neq  {k}_{2}$ 知 $m$ 和 $n$ 不平行,设 $m\text{、}n$ 相交于点 $P$，则直线 $m$ 以 $P$ 为中心旋转后与 $n$ 重合,所以 ${C}$ 正确
    \end{jieda}
    
    \item 已知直线 ${l}_{1},{l}_{2}$ 的斜率是方程 ${x}^{2} - {px} - 2 = 0$ 的两个根,则\dotfill{(\kh{C})}
    \xx{${l}_{1} \bot  {l}_{2}$}{${l}_{1}//{l}_{2}$}{${l}_{1}$ 与 ${l}_{2}$ 相交但不垂直}{${l}_{1}$ 与 ${l}_{2}$ 的位置关系不确定}
    \begin{jieda}
        设直线 ${l}_{1},{l}_{2}$ 的斜率为 ${k}_{1},{k}_{2}$ ,则 ${k}_{1}{k}_{2} =  - 2$，$\therefore {k}_{1}{k}_{2} \neq   - 1,\therefore {l}_{1},{l}_{2}$ 不垂直, ${A}$ 错误;
        
        若 ${k}_{1} = {k}_{2}$ ,则 ${k}_{1}{k}_{2} = {k}_{1}^{2} \geq  0$ ,与 ${k}_{1}{k}_{2} =  - 2$ 矛盾, $\therefore {k}_{1} \neq  {k}_{2},\therefore {l}_{1},{l}_{2}$ 不平行, ${B}$ 错误;
        
        $\because {l}_{1},{l}_{2}$ 不平行,也不垂直, $\therefore {l}_{1},{l}_{2}$ 相交但不垂直, ${C}$ 正确, ${D}$ 错误.
    \end{jieda}
    
    \item 直线 $l : \left( {2 - m}\right) x + \left( {{2m} + 1}\right) y + {3m} + 4 = 0$ 分别与 $x$ 轴, $y$ 轴交于 $A\text{ 、 }B$ 两点，若三角形 ${AOB}$ 面积为 5,则实数 $m$ 的解有$\overline{\quad\quad}$个\dotfill{(\kh{D})}
    \xx{1}{2}{3}{4}
    \begin{jieda}
        由题可知,直线 $l : \left( {2 - m}\right) x + \left( {{2m} + 1}\right) y + {3m} + 4 = 0$ 的斜率存在且不为 0，故 $\left\{  \begin{array}{l} 2 - m \neq  0 \\  {2m} + 1 \neq  0 \end{array}\right.$ ,即 $m \neq  2$ 且 $m \neq   - \dfrac{1}{2}$
        
        令 $x = 0$ ,得 $y =  - \dfrac{{3m} + 4}{{2m} + 1}$ ; 令 $y = 0$ ,得 $x = \dfrac{{3m} + 4}{m - 2}$ ; 即 $B\left( {0, - \dfrac{{3m} + 4}{{2m} + 1}}\right) ,A\left( {\dfrac{{3m} + 4}{m - 2},0}\right)$，所以 ${S}_{\bigtriangleup {AOB}} = \dfrac{1}{2}{OA} \cdot  {OB} = \dfrac{1}{2}\left| \dfrac{{3m} + 4}{m - 2}\right|  \cdot  \left| {-\dfrac{{3m} + 4}{{2m} + 1}}\right|  = 5$ ,所以 $9{m}^{2} + {24m} + {16} = {10}\left| {2{m}^{2} - {3m} - 2}\right|$，则 ${11}{m}^{2} - {54m} - {36} = 0$ 或 ${29}{m}^{2} - {6m} - 4 = 0$ ,解得 $m = \dfrac{{27} \pm  {15}\sqrt{5}}{11}$ 或 $m = \dfrac{3 \pm  5\sqrt{5}}{29}$，故解得的实数 $m$ 的解有 4 个.
    \end{jieda}
    
    \item 已知 $\bigtriangleup  {{ABC}}$ 三边所在直线方程为 ${AB} : {3x} + {4y} + {12} = 0,{BC} : {4x} - {3y} + {16} = 0$，${AC} : {2x} + y - 2 = 0.$
    \begin{enumerate}
        \item 求 ${AC}$ 边上的高所在的直线方程
        \item 求直线 ${AB}$ 与直线 ${AC}$ 的夹角
    \end{enumerate}
    \begin{jieda}
        \begin{enumerate}
            \item 
            如图所示：
            \begin{center}
                \includegraphics[width=0.3\linewidth]{bodyxb//src/zhixiandeweizhiguanxi2.png}
            \end{center}
            由 $\left\{  \begin{array}{l} {3x} + {4y} + {12} = 0 \\  {4x} - {3y} + {16} = 0 \end{array}\right.$ ,可得 $\left\{  \begin{array}{l} x =  - 4 \\  y = 0 \end{array}\right.$ ，即 $B\left( {-4,0}\right)$ ;
            又因为 ${k}_{AC} =  - 2$，所以 ${AC}$ 边上的高所在直线的斜率 $k = \dfrac{1}{2}$，所以 ${AC}$ 边上的高所在的直线方程为 $y = \dfrac{1}{2}\left( {x + 4}\right)$，即 $x - {2y} + 4 = 0$ 
            
            \item 因为 ${k}_{AC} =  - 2,{k}_{AB} =  - \dfrac{3}{4}$ 
            
            设直线 ${AB}$ 与直线 ${AC}$ 的夹角为 $\theta$ ,
            则有 $\tan \theta  = \dfrac{\left| {k}_{AC} - {k}_{AB}\right| }{1 + {k}_{AC} \cdot  {k}_{AB}} = \dfrac{1}{2}$，所以 $\theta  = \arctan \dfrac{1}{2}$ .
        \end{enumerate}
    \end{jieda}
    
    \item 等腰直角 $\bigtriangleup  {ABC}$ 的直角顶点 $C$ 和顶点 $B$ 都在直线 ${2x} + {3y} - 6 = 0$ 上,顶点 $A$ 的坐标是$(1, -2)$,求边 ${AB},{AC}$ 所在的直线方程.
    \begin{jieda}
        ${AB} : x - {5y} - {11} = 0$ 或 ${5x} + y - 3 = 0$
        
        ${AC} : {3x} - {2y} - 7 = 0$
    \end{jieda}
    
    \item 求过点 $P\left( {5, - 2}\right)$ 与直线 $x - y + 5 = 0$ 相交成 ${45}^{ \circ  }$ 角的直线 $l$ 的方程.
    \begin{jieda}
        $x = 5$或$y = -2$
    \end{jieda}
    
    \item 已知对于直线 $l$ 上的任意一点 $P\left( {x,y}\right)$ ,点 $Q\left( {{4x} + {2y},x + {3y}}\right)$ 仍在该直线上，求直线 $l$ 的方程.
    \begin{jieda}
        $x+y = 0$或$x - 2y = 0$
    \end{jieda}
    
    \item 在足球比赛中,甲队的边锋从乙队球门的一侧带球过人,沿平行于边线的直线向前推进. 边锋在何处射门可命中球门的角最大?
    
    \begin{center}
    \includegraphics[width=0.3\textwidth]{bodyxb/src/bo_d2oeegjef24c73b1lib0_3_239_1460_589_290_0.jpg}
    \end{center}
    \begin{jieda}
        如图所示：
        \begin{center}
            \includegraphics[width=0.3\linewidth]{bodyxb//src/zhixiandeweizhiguanxi3.png}
        \end{center}
        
        设球门的两立柱分别为点 $A$ 、点 $B$ ,边锋为 $C$ . 由 $C$ 引球门所在直线的垂线,垂足为 $O$ . 设 ${OC} = x,{OB} = b,{OA} = a,x > 0,a > b > 0$ ,且 $a\text{、}b$ 为定值. $\angle {ACO} = \alpha ,\angle {BCO} = \beta ,\angle {ACB} = \alpha  - \beta  = \gamma$ ,显然 $0 < \gamma  < \dfrac{\pi }{2}$ . 因此,$\tan \alpha  = \dfrac{a}{x},\tan \beta  = \dfrac{b}{x}$ . 所以, $\tan \gamma  = \tan \left( {\alpha  - \beta }\right)  = \dfrac{\tan \alpha  - \tan \beta }{1 + \tan \alpha  \cdot  \tan \beta } = \dfrac{\left( {a - b}\right) x}{{x}^{2} + {ab}} = \dfrac{a - b}{x + \dfrac{ab}{x}} \leq  \dfrac{a - b}{2\sqrt{ab}}$ ,当且仅当 $x = \sqrt{ab}$ 时,上述等式成立. 
        
        当 $0 < \gamma  < \dfrac{\pi }{2}$ 时, $f\left( \gamma \right)  = \tan\gamma$为增函数. 因此,当 $x = \sqrt{ab}$ 时, $\tan \gamma$ 达最大值,从而 $\angle {ACB}$ 达最大值$arctan \dfrac{a - b}{2\sqrt{ab}}$ .
    \end{jieda} 
\end{enumerate}

\clearpage
\section*{点到直线的距离}
\subsection*{点到直线的距离(1)}

\zsd 点到直线的距离公式

A 组：

\begin{enumerate}
    \item 点 $P_0\left( x_0, y_0\right)$ 到直线 $l : Ax + By + C = 0$ 的距离\blkda[4]{$\dfrac{|Ax_0+By_0+C|}{\sqrt{A^2+B^2}}$}.
    
    \item 点$(2,1)$到直线 ${3x} - {4y} + 3 = 0$ 的距离为\blkda{1}.

    \item 若点 $P\left( {-1,a}\right)$ 到直线 ${3x} + y - 6 = 0$ 的距离为 $\sqrt{10}$ ,则 $a$ 的值为\blkda{$-1$或$19$}.
    
    \item 点$(2,5)$到直线 ${mx} - {2y} + m + 4 = 0$ 的距离的最大值为\blkda{$3\sqrt{2}$}.
    
    \item 与直线 ${6x} - {8y} + 3 = 0$ 垂直,且与原点距离等于 1 的直线方程是\blkda[3]{$4x + 3y \pm 5 = 0$}.
    
    \item 过点 $P\left( {-2,3}\right)$ 且和原点距离是 2 的直线方程是\blkda[5]{$x = -2$或$5x+12y-26 = 0$}.
\end{enumerate}

B 组：
\begin{enumerate}
    \item 若点 $P\left( {1,1}\right)$ 到直线 $x\cos \theta  + y\sin \theta  = 2$ 的距离为 $d$ ,则 $d$ 的最大值是\blkda{$2 + \sqrt{2}$}.
    \begin{jieda}
        $d = \dfrac{|cos\theta + sin\theta - 2|}{\sqrt{cos^2\theta + sin^2\theta}}  = \left|\sqrt{2}sin\left(\theta+\dfrac{\pi}{4}\right) - 2\right| \leq 2 + \sqrt{2}$
    \end{jieda}
    
    \item 在直线 $x + {3y} = 0$ 上求一点,使它到原点的距离和到直线 $x + {3y} + 2 = 0$ 的距离相等,则此点的坐标是\blkda[4]{$\left(-\dfrac{3}{5}, \dfrac{1}{5}\right)$或$\left(\dfrac{3}{5}, -\dfrac{1}{5}\right)$}.
    \begin{jieda}
        设该点坐标为$(-3m, m)$，则根据题意有$\sqrt{(-3m)^2+m^2} = \dfrac{|-3m + 3m + 2|}{\sqrt{1+9}}$，故$|m| = \dfrac{1}{5}$，故点的坐标为$\left(-\dfrac{3}{5}, \dfrac{1}{5}\right)$或$\left(\dfrac{3}{5}, -\dfrac{1}{5}\right)$
    \end{jieda}
    
    \item 若三角形的三个顶点的坐标分别为 $A\left( {-4, - 2}\right)$ 、 $B\left( {1,3}\right)$ 、 $C\left( {-5,6}\right)$ ,则 $\bigtriangleup  {ABC}$ 的面积为\blkda{$\dfrac{45}{2}$}.
    \begin{jieda}
        $\vv{AB} = (5,5), \vv{AC} = (-1, 8)$，$S = \dfrac{1}{2} |5\times8 - 5 \times (-1)| = \dfrac{45}{2}$
    \end{jieda}
    
    \item 若 $x + y + 1 = 0$ ,则 $\sqrt{{\left( x - 1\right) }^{2} + {\left( y - 1\right) }^{2}}$ 的最小值为\blkda{$\dfrac{3\sqrt{2}}{2}$}.
    
    \item 若直线 ${l}_{1}$ 与直线 ${l}_{2} : {3x} - {4y} + 1 = 0$ 垂直,且点$(-3,4)$到直线 ${l}_{1}$ 的距离是 1 ,则直线 ${l}_{1}$ 的方程是\blkda[4]{$4x+3y\pm5 = 0$}.
    \begin{jieda}
        根据垂直关系可设$l_1 : 4x+3y + c = 0$，根据点$(-3,4)$到直线 ${l}_{1}$ 的距离是$1$有$\dfrac{|c|}{5} = 1$，因此$c = \pm 5$
    \end{jieda}
    
    \item 已知直线 $l$ 过点 $A\left( {3,2}\right)$ ,且点 $B\left( {-1,3}\right)$ 到直线 $l$ 的距离等于 4,则直线 $l$ 的倾斜角为\blkda[3]{$arctan \dfrac{15}{8}$或$\dfrac{\pi}{2}$}.
    \begin{jieda}
        \begin{dingautolist}{172}
            \item 若直线$l$的斜率不存在，则$l : x = 3$，满足要求.
            \item 若直线$l$的斜率存在，则$l : y = k(x-3) + 2$，故有$\dfrac{|-4k-3+2|}{\sqrt{1+k^2}} = 4$，故$(4k+1)^2 = 16(1+k^2)$，即$16k^2+8k+1 = 16+16k^2$，解得$k = \dfrac{15}{8}$
        \end{dingautolist}
        综上，直线 $l$ 的倾斜角为$arctan \dfrac{15}{8}$或$\dfrac{\pi}{2}$
    \end{jieda}
    
    \item 在直线 ${x} - {y} + 4 = 0$ 上取一点 ${P}$ ,使它到点 ${M}\left( {-2, - 4}\right) ,{N}\left( {4,6}\right)$ 的距离相等,则点 ${P}$ 的坐标为\blkda{$\left( {-\dfrac{3}{2},\dfrac{5}{2}}\right)$}.
    \begin{jieda}
        设直线 $x - y + 4 = 0$ 上一点 $P\left( {x,x + 4}\right)$ ,则 $P$ 到点 $M\left( {-2, - 4}\right) ,N\left( {4,6}\right)$ 的距离相等，$\therefore \sqrt{{\left( x + 2\right) }^{2} + {\left( x + 4 + 4\right) }^{2}} = \sqrt{{\left( x - 4\right) }^{2} + {\left( x + 4 - 6\right) }^{2}}$，解得 $x =  - \dfrac{3}{2},\therefore y =  - \dfrac{3}{2} + 4 = \dfrac{5}{2}$，$\therefore$ 点 $P$ 的坐标为 $\left( {-\dfrac{3}{2},\dfrac{5}{2}}\right)$
    \end{jieda}
    
    \item 已知直线 ${l}_{1} : {7x} + {5y} - {24} = 0$ 与直线 ${l}_{2} : x - y = 0$ 的交点为 $A$ ,直线 ${l}_{3}$ 过点 $A$ ,且点(5,1)到直线 ${l}_{3}$ 的距离为 3 ,则直线 ${l}_{3}$ 的方程为\blkda[4]{$4x-3y-2 = 0$或$x = 2$}.
    
    \item 已知 $A\text{ 、 }B\text{ 、 }C$ 为三角形的三个内角. 若原点到直线 $l : x\sin A + y\sin B + \sin C = 0$ 的距离大于 1 ,则此三角形为\dotfill{(\kh{C})}
    \xx{锐角三角形}{直角三角形}{钝角三角形}{等边三角形}
    
    \item 若点 $M\left( {\cos \alpha,\sin \alpha}\right)$ 、点 $N\left( {\sin \alpha,\cos \alpha}\right)$ 到 直线$x\sin \alpha  + y\cos \alpha  + p = 0\left( {p <  - 1}\right)$ 的距离分别为 $m,n$ ,则 $m$ 与 $n$ 的大小关系是\dotfill{(\kh{A})}
    \xx{$m \geq  n$}{$m \leq  n$}{$m \neq  n$}{以上结论都不对}
    \begin{jieda}
        $m=\dfrac{|\cos \alpha\sin \alpha  + \sin \alpha \cos \alpha  + p|}{\sqrt{sin^2\alpha + cos^2\alpha}} = |p+sin2\alpha|$，$n=\dfrac{|\sin \alpha \sin \alpha  + \cos \alpha\cos \alpha  + p|}{\sqrt{sin^2\alpha + cos^2\alpha}} = |p+1|$

        因为$p < 1$，所以$n \leq m$
    \end{jieda}
    
    \item 设 $M\left( {{x}_{1},{y}_{1}}\right) ,N\left( {{x}_{2},{y}_{2}}\right)$ 为不同的两点,直线 $l : {ax} + {by} + c = 0,\delta  = \dfrac{a{x}_{1} + b{y}_{1} + c}{a{x}_{2} + b{y}_{2} + c}$ , 以下命题中正确的序号为\dotfill{(\kh{D})}
    \begin{dingautolist}{172}
        \item 存在实数 $\delta$ ,使得点 $N$ 在直线 $l$ 上
        \item 若 $\delta  = 1$ ,则过 $M$ 、 $N$ 的直线与直线 $l$ 平行
         \item 若 $\delta  =  - 1$ ,则直线 $l$ 经过 ${MN}$ 的中点
        \item 若 $0 < \delta  < 1$ ,则点 $M$ 、 $N$ 在直线 $l$ 的同侧且直线 $l$ 与线段 ${MN}$ 的反向延长线相交
    \end{dingautolist}
    \xx{\ding{172}\ding{173}\ding{174}}{\ding{173}\ding{175}}{\ding{174}\ding{175}}{\ding{173}\ding{174}\ding{175}}

    
    \begin{jieda}
        \begin{dingautolist}{172}
            \item 点 $N$ 在直线 $l$ 上$\Longleftrightarrow $分母为0，故不可能.
            \item $\delta  = 1 \Longleftrightarrow a(x_1-x_2)+b(y_1-y_2) = 0$，故过 $M,N$ 两点的直线与直线 $l$ 平行（法向量相同），但点 $N$ 不在直线 $l$ 上，故必然是平行关系.
            \item $\delta  = -1 \Longleftrightarrow a(x_1+x_2)+b(y_1+y_2) + 2c = 0 \Longleftrightarrow a \cdot \dfrac{x_1+x_2}{2} + b \cdot \dfrac{y_1+y_2}{2} + c = 0$，故直线 $l$ 经过线段 ${MN}$ 的中点.
            \item 当 $0 < \delta  < 1$ ,得到 $0 < \dfrac{a{x}_{1} + b{y}_{1} + c}{a{x}_{2} + b{y}_{2} + c} < 1$ ,即 $\left( {a{x}_{1} + b{y}_{1} + c}\right) \left( {a{x}_{2} + b{y}_{2} + c}\right)  > 0$ ,所以得到点 $M, N$ 在直线 $l$ 的同侧,且 $\left| {a{x}_{1} + b{y}_{1} + c}\right|  < \left| {a{x}_{2} + b{y}_{2} + c}\right|$ ,得到点 $M$ 与点 $N$ 到直线 $l$ 的距离不等，所以直线 $l$ 与线段 ${MN}$ 的反向延长线相交
        \end{dingautolist}
    \end{jieda}
    
    \item 已知直线 $l$ 的倾斜角为 $a,\cos \alpha  = \dfrac{\sqrt{2}}{2}$ ,且这条直线经过点 $A\left( {\sqrt{3},2}\right)$ .
    \begin{enumerate}
        \item 求直线 $l$ 的方程
        \item 直线 ${kx} - y + 1 - \sqrt{3}k = 0$ 恒过定点 $B$ ,求点 $B$ 到直线 $l$ 的距离
    \end{enumerate}
    \begin{jieda}
        \begin{enumerate}
            \item 由 $\cos \alpha  = \dfrac{\sqrt{2}}{2},a \in  \lbrack 0,\pi )$ ,则 $a = \dfrac{\pi }{4},\tan \alpha  = 1$，$\therefore$ 直线 $l$ 的斜率 $k = 1$ ,且直线过点 $A\left( {\sqrt{3},2}\right)$，$\therefore$ 由直线的点斜式方程得 $l : y - 2 = \left( {x - \sqrt{3}}\right)$，即 $x - y + 2 - \sqrt{3} = 0$，$\therefore$ 所求直线 $l$ 的方程为 $x - y + 2 - \sqrt{3} = 0$
            
            \item $\because$ 直线 ${kx} - y + 1 - \sqrt{3}k = 0$ 化简得： $k\left( {x - \sqrt{3}}\right)  - y + 1 = 0$，$\therefore$ 定点 $B\left( {\sqrt{3},1}\right)$，则点 $B\left( {\sqrt{3},1}\right)$ 到直线 $x - y + 2 - \sqrt{3} = 0$ 的距离为:$d = \dfrac{\left| \sqrt{3} \times  1 - 1 + 2 - \sqrt{3}\right| }{\sqrt{{1}^{2} + {1}^{2}}} = \dfrac{1}{\sqrt{2}} = \dfrac{\sqrt{2}}{2}$，故 $B$ 到直线 $l$ 的距离为 $\dfrac{\sqrt{2}}{2}$ .
        \end{enumerate}
    \end{jieda}
    
    \item 已知 $A\left( {1,3}\right) ,B\left( {5,7}\right)$
    \begin{enumerate}
        \item 求线段 ${AB}$ 垂直平分线所在直线方程
        \item 若直线 $l$ 过 $(-1,0)$,且 ${A}$ 、 $B$ 到直线 $l$ 距离相等,求 $l$ 方程
    \end{enumerate}
    \begin{jieda}
        \begin{enumerate}
            \item 因为点 $A\left( {1,3}\right)$ , $B\left( {5,7}\right)$ .所以线段 ${AB}$ 的中点坐标为$(3,5)$,直线 ${AB}$ 的斜率为 $\frac{7 - 3}{5 - 1} = 1$，因此直线 ${AB}$ 的中垂线的斜率为 $-1$，因此线段 ${AB}$ 的垂直平分线所在直线方程为 $y - 5 =  - \left( {x - 3}\right)$，即 $x + y - 8 = 0$
            \item 因为直线 $l$ 过点$(-1,0)$, $A\left( {1,3}\right)$ , $B\left( {5,7}\right)$，当直线 $l$ 的斜率不存在时,显然不合题意
            
            设直线 $l$ 的方程为 $y = k\left( {x + 1}\right)$ ,即 ${kx} - y + k = 0$，所以 $\dfrac{\left| k - 3 + k\right| }{\sqrt{{k}^{2} + {\left( -1\right) }^{2}}} = \dfrac{\left| 5k - 7 + k\right| }{\sqrt{{k}^{2} + {\left( -1\right) }^{2}}}$ ,解得 $k = 1$ 或 $k = \dfrac{5}{4}$，所以直线 $l$ 的方程为 $x - y + 1 = 0$ 或 ${5x} - {4y} + 5 = 0$ .
        \end{enumerate}
    \end{jieda}
    \item 已知 $\bigtriangleup  {ABC}$ 的三个顶点 $A\left( {m,n}\right)$ 、 $B\left( {2,1}\right)$ 、 $C\left( {-2,3}\right)$ .
    \begin{enumerate}
        \item 求 ${BC}$ 边所在直线的方程
        \item ${BC}$ 边上中线 ${AD}$ 的方程为 ${2x} - {3y} + 6 = 0$ ,且 ${S}_{\bigtriangleup {ABC}} = 7$ ,求点 $A$ 的坐标
    \end{enumerate}
    \begin{jieda}
        \begin{enumerate}
            \item 因为 $B\left( {2,1}\right) \text{、}C\left( {-2,3}\right)$ ,所以 ${BC}$ 边所在直线的方程为：$\dfrac{y - 1}{3 - 1} = \dfrac{x - 2}{-2 - 2} \Longrightarrow  x + {2y} - 4 = 0$
            
            \item ${BC}$ 边上中线 ${AD}$ 的方程为 ${2x} - {3y} + 6 = 0$ ,所以有 ${2m} - {3n} + 6 = 0$ 

            $\vv{CB} = (4, -2), \vv{CA} = (m+2, n-3)$，$S = \dfrac{1}{2} \times |4(n-3) - (-2)(m+2)| = 7 \Longrightarrow \left| {m + {2n} - 4}\right|  = 7$
            
            因此有 $\left\{  \begin{array}{l} m + {2n} - 4 = 7 \\  {2m} - {3n} + 6 = 0 \end{array}\right.$ 或 $\left\{  \begin{array}{l} m + {2n} - 4 =  - 7 \\  {2m} - {3n} + 6 = 0 \end{array}\right.$ ,解得： $\left\{  \begin{array}{l} m = 3 \\  n = 4 \end{array}\right.$ 或 $\left\{  \begin{array}{l} m =  - 3 \\  n = 0 \end{array}\right.$，所以点 $A$ 的坐标为:$(3,4)$或$(-3,0)$.
        \end{enumerate}
    \end{jieda}
    
    \item 过点 $P\left( {2,1}\right)$ 的直线 $l$ 分别交 $y = x\left( {x \geq  0}\right)$ 与 $y =  - x\left( {x \geq  0}\right)$ 于 $A\text{ 、 }B$ 两点.
    \begin{enumerate}
        \item 若直线 $l$ 的倾斜角为 $\dfrac{\pi }{4}$ ,求直线 $l$ 的一般式方程.
        \item 当 $\left| {PA}\right|  \cdot  \left| {PB}\right|$ 最小时,求直线 $l$ 的方程
        \item 已知 $O$ 为坐标原点,设 $\triangle{AOB}$ 的面积为 ${S}$ ,讨论这样的直线 $l$ 的条数
    \end{enumerate}
    \begin{jieda}
        \begin{enumerate}
            \item 若直线 $l$ 的倾斜角为 $\dfrac{\pi }{4}$ ,则直线 $l$ 的方程为 $y - 1 = 1\left( {x - 2}\right)$ ,即 $x - y - 1 = 0$ 
            
            \item 
            \begin{enumerate}[1.]
                \item 方法1 
                \begin{dingautolist}{172}
                    \item 当直线 $l$ 的斜率不存在时, $\left| {PA}\right| \left| {PB}\right|  = 3$
                    \item 当直线 $l$ 的斜率存在时,设直线 $l : y - 1 = k\left( {x - 2}\right) ,k \in  \left( {-\infty , - 1}\right)  \cup  \left( {1, + \infty }\right)$ ,
                    $\left\{  {\begin{matrix} y = x \\  y - 1 = k\left( {x - 2}\right)  \end{matrix}\text{ 得 }A\left( {\dfrac{{2k} - 1}{k - 1},\dfrac{{2k} - 1}{k - 1}}\right) ,\left\{  {\begin{matrix} y =  - x \\  y - 1 = k\left( {x - 2}\right)  \end{matrix}\text{ 得 }B\left( {\dfrac{{2k} - 1}{k + 1},\dfrac{1 - {2k}}{k + 1}}\right) ,}\right. }\right.$
    
                    $\vv{PA} = \left( {\dfrac{1}{k - 1},\dfrac{k}{k - 1}}\right) ,\vv{PB} = \left( {\dfrac{-3}{k + 1},\dfrac{-{3k}}{k + 1}}\right)$ ,且 $\vv{PA},\vv{PB}$ 反向，所以 $\left| {PA}\right| \left| {PB}\right|  = \left| {\vv{PA} \cdot  \vv{PB}}\right|  = \dfrac{3 + 3{k}^{2}}{\left| {k}^{2} - 1\right| } = 3 + \dfrac{6}{\left| {k}^{2} - 1\right| } > 3$ 
                \end{dingautolist}
                综上所述: $\left| {PA}\right|  \cdot  \left| {PB}\right|$ 的最小值为 3,此时直线 $l$ 的斜率不存在,直线方程为 $x = 2$ .
                \item 方法2
                
                设直线$l$与$y = x$的夹角为$\theta$，$\theta \in \left(0, \dfrac{\pi}{2}\right)$则$|PA| = \dfrac{1}{\sqrt{2}} \times \dfrac{1}{sin\theta}$，$|PB| = \dfrac{3}{\sqrt{2}} \times \dfrac{1}{cos\theta}$，故$|PA|\cdot|PB| = \dfrac{3}{sin2\theta} \leq 3$，当$\theta = \dfrac{\pi}{4}$时取得最小值$3$，此时直线方程为$x = 2$.
            \end{enumerate}
            
            \item 
            \begin{dingautolist}{172}
                \item 当直线斜率不存在时, $A\left( {2,2}\right) ,B\left( {2, - 2}\right) ,S = \dfrac{1}{2}\left| {OA}\right|  \cdot  \left| {OB}\right|  = 4$
                \item 当直线斜率存在时, $k \in  \left( {-\infty , - 1}\right)  \cup  \left( {1, + \infty }\right) ,A\left( {\dfrac{{2k} - 1}{k - 1},\dfrac{{2k} - 1}{k - 1}}\right) ,B\left( {\dfrac{{2k} - 1}{k + 1},\dfrac{1 - {2k}}{k + 1}}\right)$ ,
                
                $S = \dfrac{1}{2} \left|\dfrac{{2k} - 1}{k - 1} \cdot \dfrac{1 - {2k}}{k + 1} - \dfrac{{2k} - 1}{k - 1} \cdot \dfrac{{2k} - 1}{k + 1}\right| = \dfrac{{\left( 2k - 1\right) }^{2}}{{k}^{2} - 1}$，即 $\left( {4 - S}\right) {k}^{2} - {4k} + S + 1 = 0$
                \begin{enumerate}[1.]
                    \item 方法1 
                    
                    当 $S = 4$ 时,方程有 1 解,此时 $k = \dfrac{5}{4}$
                    
                    当 $S \neq  4$ 时, $\Delta  = {16} - 4\left( {4 - S}\right) \left( {S + 1}\right)  = {4S}\left( {S - 3}\right)$ 
                    
                    当 $S < 3$ 时, $\Delta  < 0$ ,方程无解
                    
                    当 $S = 3$ 时, $\Delta  = 0,k = 2$ ,方程有 1 解
                    
                    当 $4 > S > 3$ 时, $\Delta  > 0,f\left( k\right)  = \left( {4 - S}\right) {k}^{2} - {4k} + S + 1$，开口向上，对称轴 $\dfrac{2}{4 - S} > 2$ ,且 $f\left( 1\right)  = 1 > 0$ ,方程有两个大于 1 的解.
                    
                    当 $S > 4$ 时, $\Delta  > 0,f\left( k\right)  = \left( {4 - S}\right) {k}^{2} - {4k} + S + 1$ 开口向下, $f\left( 1\right)  = 1 > 0,f\left( {-1}\right)  = 9 > 0$ ,方程有 1 个大于 1 的解, 一个小于 -1 的解.
                    \item 方法2

                    $S = \dfrac{(2k-1)^2}{k^2-1}$
                    
                    令$t = 2k-1$，$t \in (-\infty, -3) \cup (1, +\infty)$，$t$与$k$一一对应，则$S = \dfrac{t^2}{\left(\dfrac{t+1}{2}\right)^2-1} = \dfrac{4t^2}{t^2+2t-3} = \dfrac{4}{1+\dfrac{2}{t}-\dfrac{3}{t^2}}$，即$1+\dfrac{2}{t}-\dfrac{3}{t^2} = \dfrac{4}{S}$

                    令$m = \dfrac{1}{t}$，$m \in \left(-\dfrac{1}{3}, 0\right) \cup (0, 1)$，$t$与$m$一一对应，则$-3m^2+2m+1 = \dfrac{4}{S}$
                    
                    $g(m) = -3m^2+2m+1$在$\left(-\infty, \dfrac{1}{3}\right)$严格增，在$\left(\dfrac{1}{3}, +\infty\right)$严格减
                    
                    数形结合可知当$\dfrac{4}{S} \in [0, 1)\cup \left(1, \dfrac{4}{3}\right)$时原方程有两解，当$\dfrac{4}{S} \in \left\{1, \dfrac{4}{3}\right\}$时，原方程只有一解，即$S \in (3, 4) \cup (4, +\infty)$时原方程有两解，当$S \in \{3,4\}$原方程只有一解.
                \end{enumerate}
            \end{dingautolist}

            综上所述：当 $S < 3$ 时, 0 条；当 $S = 3$ 时, 1 条；当 $S > 3$ 时, 2 条.
        \end{enumerate}
    \end{jieda}
\end{enumerate}

C 组：
\begin{enumerate}
    \item 已知矩形 ${OABC}$ 中, $O$ 为坐标原点,点 $A$ 在 $x$ 轴上,点 $C$ 在 $y$ 轴上, $B$ 的坐标为$(10,5)$,点 $P$ 在边 ${BC}$ 上,点 $A$ 关于 ${OP}$ 的对称点为 ${A}^{\prime }$ ,若点 ${A}^{\prime }$ 到直线 ${BC}$ 的距离为 4 ,则点 ${A}^{\prime }$ 的坐标可能为\blkda[6]{$\left( {-3\sqrt{11},1}\right) ,\left( {-\sqrt{19},9}\right) ,\left( {\sqrt{19},9}\right)$}.
    \begin{jieda}
    依题意,点 $A\left( {{10},0}\right)$ ,直线 ${BC} : y = 5$ ,而点 $P$ 在边 ${BC}$ 上,则直线 ${OP}$ 的斜率 $k \geq  \dfrac{5}{10} = \dfrac{1}{2}$或 ${OP}$ 在 $y$ 轴上，设点 ${A}^{\prime }\left( {a,b}\right)$ ,由点 ${A}^{\prime }$ 到直线 ${BC}$ 的距离为 4,得 $\left| {b - 5}\right|  = 4$ ,即 $b = 1$ 或 $b = 9$
    
    又点 $A$ 关于 ${OP}$ 的对称点为 ${A}^{\prime }$ ,则 $\left| {O{A}^{\prime }}\right|  = \left| {OA}\right|  = {10}$ ,即 ${a}^{2} + {b}^{2} = {10}^{2}$ 

    \begin{dingautolist}{172}
        \item 当 $b = 1$ 时, $a =  - 3\sqrt{11}$ 或 $a = 3\sqrt{11}$
        
        若 $a =  - 3\sqrt{11}$ ,有 ${A}^{\prime }\left( {-3\sqrt{11},1}\right)$ ,点 ${A}^{\prime }$ 与 $A$ 的中点$\left( {\dfrac{{10} - 3\sqrt{11}}{2},\dfrac{1}{2}}\right)$ 在直线 ${OP}$ 上，此时直线 ${OP}$ 斜率 $k = \dfrac{1}{{10} - 3\sqrt{11}} = {10} + 3\sqrt{11} > \dfrac{1}{2}$ ,符合题意,则 ${A}^{\prime }\left( {-3\sqrt{11},1}\right)$
        
        若 $a = 3\sqrt{11}$ ,有 ${A}^{\prime }\left( {3\sqrt{11},1}\right)$ ,点 ${A}^{\prime }$ 与 $A$ 的中点 $\left( {\dfrac{{10} + 3\sqrt{11}}{2},\dfrac{1}{2}}\right)$ 在直线 ${OP}$ 上，此时直线 ${OP}$ 斜率 $k = \dfrac{1}{{10} + 3\sqrt{11}} = {10} - 3\sqrt{11} < \dfrac{1}{2}$ ,不符合题意
        \item 当 $b = 9$ 时, $a =  - \sqrt{19}$ 或 $a = \sqrt{19}$
        
        若 $a =  - \sqrt{19}$ ,有 ${A}^{\prime }\left( {-\sqrt{19},9}\right)$ ,点 ${A}^{\prime }$ 与 $A$ 的中点 $\left( {\dfrac{{10} - \sqrt{19}}{2},\dfrac{9}{2}}\right)$在直线 ${OP}$ 上.此时直线 ${OP}$ 斜率 $k = \dfrac{9}{{10} - \sqrt{19}} = \dfrac{{10} + \sqrt{19}}{9} > \dfrac{1}{2}$ ,符合题意,则 ${A}^{\prime }\left( {-\sqrt{19},9}\right)$ ,
        
        若 $a = \sqrt{19}$ ,有 ${A}^{\prime }\left( {\sqrt{19},9}\right)$ ,点 ${A}^{\prime }$ 与 $A$ 的中点 $\left( {\dfrac{{10} + \sqrt{19}}{2},\dfrac{9}{2}}\right)$ 在直线 ${OP}$ 上，此时直线 ${OP}$ 斜率 $k = \dfrac{9}{{10} + 3\sqrt{11}} = \dfrac{{10} - \sqrt{19}}{9} > \dfrac{1}{2}$ ,符合题意,则 ${A}^{\prime }\left( {\sqrt{19},9}\right)$
    \end{dingautolist}
    所以点 ${A}^{\prime }$ 的坐标可能为 $\left( {-3\sqrt{11},1}\right) ,\left( {-\sqrt{19},9}\right) ,\left( {\sqrt{19},9}\right)$ .
    \end{jieda}

    \item 我国著名数学家华罗庚曾说“数缺形时少直观,形少数时难入微；数形结合百般好, 隔离分家万事休. ” 事实上,很多代数问题可以都转化为几何问题加以解决,例如,与 $\sqrt{{\left( x - a\right) }^{2} + {\left( y - b\right) }^{2}}$ 相关的代数问题. 可以转化为点$(x, y)$与点$(a, b)$之间的距离的几何问题. 已知点 $M\left( {{x}_{1},{y}_{1}}\right)$ 在直线 ${l}_{1} : y = x + 2$ ,点 $N\left( {{x}_{2},{y}_{2}}\right)$ 在直线 ${l}_{2} : y = x$ 上,且 ${MN} \bot  {l}_{1}$ ,结合上述观点, $\sqrt{{x}_{1}^{2} + {\left( {y}_{1} - 4\right) }^{2}} + \sqrt{{\left( {x}_{2} - 5\right) }^{2} + {y}_{2}^{2}}$ 的最小值为\blkda{5}.
    \begin{jieda}
        $\sqrt{{x}_{1}^{2} + {\left( {y}_{1} - 4\right) }^{2}}$ 表示点 $M\left( {{x}_{1},{y}_{1}}\right)$ 到点 $A\left( {0,4}\right)$ 的距离，$\sqrt{{\left( {x}_{2} - 5\right) }^{2} + {y}_{2}^{2}}$ 表示点 $N\left( {{x}_{2},{y}_{2}}\right)$ 到点 $B\left( {5,0}\right)$ 的距离，所以 $\sqrt{{x}_{1}^{2} + {\left( {y}_{1} - 4\right) }^{2}} + \sqrt{{\left( {x}_{2} - 5\right) }^{2} + {y}_{2}^{2}} = \left| {MA}\right|  + \left| {NB}\right|$
        
        \begin{center}
        \includegraphics[width=0.3\textwidth]{bodyxb/src/bo_d2ohv9f7aajc738rr47g_8_242_652_401_302_0.jpg}
        \end{center}
        
        过点 ${A}$ 作 ${AC} \bot  {l}_{1}$ ,垂足为 $C$，下证四边形 ${AMNC}$ 为平行四边形：
        
        因为直线 ${l}_{1}$ 的方程为 $x - y + 2 = 0, A\left( {0,4}\right)$，所以 $\left| {AC}\right|  = \dfrac{\left| 0 - 4 + 2\right| }{\sqrt{1 + 1}} = \sqrt{2}$
        
        又直线 ${l}_{1} : y = x + 2$ 与直线 ${l}_{2} : y = x$ 平行, ${MN} \bot  {l}_{1}$，所以 $\left| {MN}\right|  = \dfrac{\left| 2 - 0\right| }{\sqrt{1 + 1}} = \sqrt{2}$
        
        所以 ${MN}//{AC},\left| {MN}\right|  = \left| {AC}\right|$，因此四边形 ${AMNC}$ 为平行四边形
        
        
        故$\left| {AM}\right|  = \left| {CN}\right|$，所以 $\sqrt{{x}_{1}^{2} + {\left( {y}_{1} - 4\right) }^{2}} + \sqrt{{\left( {x}_{2} - 5\right) }^{2} + {y}_{2}^{2}} = \left| {CN}\right|  + \left| {NB}\right|$
        
        $\left| {CN}\right|  + \left| {NB}\right|  \geq  \left| {CB}\right|$，当且仅当 $C,N,B$ 三点共线时等号成立，所以当点 $N$ 为线段 ${CB}$ 与直线 ${l}_{2}$ 的交点时$\sqrt{{x}_{1}^{2} + {\left( {y}_{1} - 4\right) }^{2}} + \sqrt{{\left( {x}_{2} - 5\right) }^{2} + {y}_{2}^{2}}$ 取最小值,最小值为 $\left| {CB}\right|$
        
        因为过点 $A\left( {0,4}\right)$ 与直线 ${l}_{1}$ 垂直的直线的方程为 $y =  - x + 4$
        
        联立 $\left\{  \begin{array}{l} y =  - x + 4 \\  y = x + 2 \end{array}\right.$ ,可得 $\left\{  \begin{array}{l} x = 1 \\  y = 3 \end{array}\right.$，所以点 $C$ 的坐标为(1,3),所以 $\left| {CB}\right|  = \sqrt{{\left( 5 - 1\right) }^{2} + {\left( 0 - 3\right) }^{2}}$
        
        所以 $\sqrt{{x}_{1}^{2} + {\left( {y}_{1} - 4\right) }^{2}} + \sqrt{{\left( {x}_{2} - 5\right) }^{2} + {y}_{2}^{2}}$ 的最小值为5.
        
    \end{jieda}

    \item 在平面直角坐标系中,定义 $d\left( {P,Q}\right)  = \left| {{x}_{1} - {x}_{2}}\right|  + \left| {{y}_{1} - {y}_{2}}\right|$ 为 $P\left( {{x}_{1},{y}_{1}}\right) ,Q\left( {{x}_{2},{y}_{2}}\right)$ 两点之间的 “折线距离”. 已知 $O\left( {0,0}\right)$ , $A\left( {2,0}\right)$ , $B\left( {2,1}\right)$ , $C\left( {0,1}\right)$ ,点 $M$ 在矩形 ${OABC}$ 内 (含边界) 且到点 $O,B$ 的 “折线距离” 相等,则点 $M$ 的轨迹长度为\blkda{$\sqrt{2}$}.
    \begin{jieda}
        设 $M\left( {x,y}\right)$ ,因为点 $M$ 在矩形 ${OABC}$ 内 (含边界)，则 $0 \leq  x \leq  2,0 \leq  y \leq  1$.
        
        \begin{center}
        \includegraphics[width=0.3\textwidth]{bodyxb/src/bo_d2ohv9f7aajc738rr47g_8_240_1644_369_267_0.jpg}
        \end{center}
        
        因为点 $M$ 到点 $O$ , $B$ 的“折线距离”相等，所以 $\left| {x - 0}\right|  + \left| {y - 0}\right|  = \left| {x - 2}\right|  + \left| {y - 1}\right|$ ,即 $x + y = 2 - x + 1 - y$，即$x + y = \dfrac{3}{2}$
        
        当 $y = 0$ 时, $x = \dfrac{3}{2}$，当 $y = 1$ 时, $x = \dfrac{1}{2}$
        
        设 $D\left( {\dfrac{3}{2},0}\right) ,E\left( {\dfrac{1}{2},1}\right)$ ,则点 $M$ 的轨迹为线段 ${DE}$，故点 $M$ 的轨迹长度为 $\left| {DE}\right|  = \sqrt{2}$ .
    \end{jieda}
\end{enumerate}

\clearpage
\subsection*{点到直线的距离(2)}

\zsd 两平行线间距离

A组：
\begin{enumerate}
    \item 已知直线 $l$ 过点$(-2,0)$,倾斜角为 ${60}^{ \circ  }$，则坐标原点到直线 $l$ 的距离为\blkda{$\sqrt{3}$}.
    
    \item 若直线 $x + y - 2 = 0$ 和直线 ${mx} + {2y} + 9 = 0$ 平行,则 $m$ 的值为\blkda{$2$}.
    
    \item 两平行直线 ${l}_{1} : {3x} + {4y} - 2 = 0$ 与 ${l}_{2} : {6x} + 8y - 5 = 0$ 之间的距离是\blkda{$\dfrac{1}{10}$}.
    \begin{jieda}
        由两平行直线 ${l}_{1} : {3x} + {4y} - 2 = 0$ 与 ${l}_{2} : {6x} + {my} - 5 = 0$ 可得$3 \times  m - 4 \times  6 = 0$ 即 $m = 8$ 
        
        直线 ${l}_{2}$ 的方程可化为 ${3x} + {4y} - \dfrac{5}{2} = 0$，故两平行直线 ${l}_{1},{l}_{2}$ 之间的距离 $d = \dfrac{\left| -2 - \left( -\dfrac{5}{2}\right) \right| }{\sqrt{{3}^{2} + {4}^{2}}} = \dfrac{1}{10}$ .
    \end{jieda}

    \item 过点 $P\left( {0, -1}\right)$ 的直线 $l$ 与以 $A\left( {3,2}\right) ,B\left( {2, -3}\right)$ 为端点的线段 ${AB}$ 有交点,求直线 $l$ 的斜率与倾斜角的取值范围.
    \begin{jieda}
        斜率$k \in [-1, 1]$，倾斜角$\theta \in \left[0, \dfrac{\pi}{4}\right] \cup \left[\dfrac{3\pi}{4}, \pi\right)$
    \end{jieda}
\end{enumerate}

B 组：

\begin{enumerate}
    \item 已知直线 ${2x} + y - 3 = 0$ 与直线 ${4x} - {my} - 3 = 0$ 平行,则它们之间的距离是\blkda{$\dfrac{3\sqrt{5}}{10}$}.
    \begin{jieda}
        直线 ${2x} + y - 3 = 0$ 与直线 ${4x} - {my} - 3 = 0$ 平行,则 $\dfrac{4}{2} = \dfrac{-m}{1} \neq  1$ ,解得 $m =  - 2$，直线 ${4x} - {my} - 3 = 0$ 为 ${2x} + y - \dfrac{3}{2} = 0$，所以它们之间的距离是 $\dfrac{\left| -3 - \left( -\dfrac{3}{2}\right) \right| }{\sqrt{{2}^{2} + {1}^{2}}} = \dfrac{3\sqrt{5}}{10}$ .
    \end{jieda}

    \item 设 $a \in  \mathbf{R}$ ,若直线 ${2x} + y - 3 = 0$ 与直线 ${2x} + y + a = 0$ 之间的距离为 $\sqrt{5}$ ,则 $a$ 的值为\blkda{$2$或$-8$}.
    \begin{jieda}
        $d = \dfrac{|a+3|}{\sqrt{4+1}} = \sqrt{5}$，解得$a = 2$或$-8$
    \end{jieda}

    \item 平面上有四条直线,它们的方程分别是 $y = {2x} + 1,{4x} - {2y} + 5 = 0,y =  - x + 1$ , $x + y - 4 = 0$ . 则由这四条直线围成四边形的面积是\blkda{$\dfrac{3}{2}$}.
    \begin{jieda}
        因为 ${4x} - {2y} + 5 = 0 \Leftrightarrow  y = {2x} + \dfrac{5}{2}$，所以直线 $y = {2x} + 1,{4x} - {2y} + 5 = 0$ 互相平行
        
        $y =  - x + 1,x + y - 4 = 0$ 互相平行,所以四条直线围成四边形为平行四边形，记 $y = {2x} + 1$ 与直线 $y =  - x + 1,x + y - 4 = 0$ 分别交于点 $A,B$，则由 $\left\{  \begin{array}{l} y =  - x + 1 \\  y = {2x} + 1 \end{array}\right.$ 解得 $A\left( {0,1}\right)$ ,由 $\left\{  \begin{array}{l} y = {2x} + 1 \\  x + y - 4 = 0 \end{array}\right.$ 解得 $B\left( {1,3}\right)$
        \begin{center}
        \includegraphics[width=0.3\textwidth]{bodyxb/src/bo_d2ohv9f7aajc738rr47g_11_242_1003_403_336_0.jpg}
        \end{center}
        \begin{enumerate}
            \item 方法1 \\
            $\left| {AB}\right|  = \sqrt{{\left( 0 - 1\right) }^{2} + {\left( 1 - 3\right) }^{2}} = \sqrt{5}$，平行直线 $y = {2x} + 1,{4x} - {2y} + 5 = 0$间的距离为$\left|\dfrac{\dfrac{5}{2} - 1}{\sqrt{5}}\right| = \dfrac{3\sqrt{5}}{10}$，所以四边形面积为 $\sqrt{5} \times  \dfrac{3\sqrt{5}}{10} = \dfrac{3}{2}$ .
            \item 方法2 \\
            $y= -x+1$与$4x-2y+5 = 0$的交点$C\left(-\dfrac{1}{2}, \dfrac{3}{2}\right)$
            
            $\vv{AC} = \left(-\dfrac{1}{2}, \dfrac{1}{2}\right), \vv{AB} = (1, 2)$，故$S = \left|-\dfrac{1}{2} \times 2 - \dfrac{1}{2} \times 1\right| = \dfrac{3}{2}$
        \end{enumerate}
    \end{jieda}
    
    \item 已知直线 ${l}_{1} : {3x} + y = 0$ 和直线 ${l}_{2} : {3x} - y - 2 = 0$ . 若直线 ${l}_{1}$ 上的点 $A$ 到原点的距离和它到直线 ${l}_{2}$ 的距离相等,则点 $A$ 的坐标是\blkda[4]{$\left(-\dfrac{1}{2}, \dfrac{3}{2}\right)$或$\left(\dfrac{1}{8}, -\dfrac{3}{8}\right)$}.
    
    \begin{jieda}
        设点$A$坐标为$(m, -3m)$，则根据题意有$\sqrt{m^2+(-3m)^2} = \dfrac{|3m + 3m - 2|}{\sqrt{1+9}}$，故$10|m| = |6m-2|$，故点的坐标为$\left(-\dfrac{1}{2}, \dfrac{3}{2}\right)$或$\left(\dfrac{1}{8}, -\dfrac{3}{8}\right)$
    \end{jieda}

    \item 若直线 ${l}_{1}$ 与直线 ${l}_{2} : x - {2y} + 4 = 0$ 平行,且直线 ${l}_{1}$ 到直线 ${l}_{2}$ 的距离与它到直线 ${l}_{2} : x - {2y} - 6 = 0$ 的距离相等,则直线 ${l}_{1}$ 的方程是\blkda[3]{$x-2y-1=0$}.
    \begin{jieda}
        根据题意可设${l}_{1}: x-2y+c=0$，且有$\dfrac{|c-(-6)|}{\sqrt{5}} = \dfrac{|c-4|}{\sqrt{5}}$，解得$c = -1$
    \end{jieda}

    \item 平行于直线 $x - y - 2 = 0$ ,且与它的距离等于 $\sqrt{2}$ 的直线方程是\blkda[5]{$x-y=0$或$x-y-4=0$}.

    \item 已知直线 ${l}_{1}$ 过点(1,0),直线 ${l}_{2}$ 过点(3,4),且 ${l}_{1}//{l}_{2}$ . 若直线 ${l}_{1}$ 与直线 ${l}_{2}$ 的距离是 2, 则直线 ${l}_{1}$ 的方程为\blkda[4]{$x = 1$或$3x-4y-3=0$},直线 ${l}_{2}$ 的方程为\blkda[4]{$x = 3$或$3x-4y+7=0$}.
    \begin{jieda}
        \begin{dingautolist}{172}
            \item 当直线斜率不存在的时候，$l_1 : x = 1, l_2: x = 3$，满足要求.
            \item 当直线斜率存在的时候，设$l_1: y = k(x-1), l_2: y = k(x-3)+4$，则有$\dfrac{|4-2k|}{\sqrt{1+k^2}} = 2$，解得$k = \dfrac{3}{4}$，此时$l_1: 3x-4y-3=0$，$l_2: 3x-4y+7=0$
        \end{dingautolist}
    \end{jieda}

    \item 若 ${2x} + y + 3 = 0$ ,则 $\sqrt{{x}^{2} + {y}^{2} - {4x} - {6y} + {13}}$ 的最小值为\blkda{$2\sqrt{5}$}.
    \begin{jieda}
        设$A(x, y)$，则$A$在直线${2x} + y + 3 = 0$上，$\sqrt{{x}^{2} + {y}^{2} - {4x} - {6y} + {13}} = \sqrt{(x-2)^2+(y-3)^2}$，即$A$到点$(2, 3)$的距离，故最小值为点$(2, 3)$到直线${2x} + y + 3 = 0$的距离，为$2\sqrt{5}$
    \end{jieda}

    \item 若点$(-1,2)$与坐标原点在直线 ${7x} + {ay} - {24} = 0$ 的同侧,则实数 $a$ 的取值范围是\blkda{$\left(-\infty, \dfrac{31}{2}\right)$}.
    \begin{jieda}
        $(7 \times (-1) + 2a - {24})(7 \times (0) + a\times(0) - {24}) > 0$，即$-7 + 2a - {24} < 0$，故$a \in \left(-\infty, \dfrac{31}{2}\right)$
    \end{jieda}

    \item 已知点 $A\left( {8,b}\right)$ 在两平行直线 ${3x} - {4y} + 5 = 0$ 和 ${6x} - {8y} + 1 = 0$ 之间,则 $b$ 应取的整数值为\blkda{7}.
    \begin{jieda}
        ${6x} - {8y} + 1 = 0$即$3x-4y+\dfrac{1}{2} = 0$，故$(24-4b+5)\left(24-4b+\dfrac{1}{2}\right) < 0$，故$b \in (6.125, 7.25)$，又因为$b \in \mathbf{Z}$，所以$b = 7$.
    \end{jieda}

    \item 已知直线 $l$ 过点 $P\left( {1,2}\right)$ ,且被两平行直线 ${l}_{1} : {4x} + {3y} + 1 = 0$ 与 ${l}_{2} : {4x} + {3y} + 6 = 0$ 截得的线段长 $\left| {AB}\right|  = \sqrt{2}$ ,则直线 $l$ 的方程为\blkda[6]{$7x-y-5=0$或$x+7y-15=0$}.
    \begin{jieda}
        易知$l_1, l_2$的距离为$1$，又因为截得的线段长 $\left| {AB}\right|  = \sqrt{2}$，所以$l$与$l_1$的夹角为$\dfrac{\pi}{4}$，故$k = \dfrac{-\dfrac{4}{3} \pm 1}{1 \mp \left(-\dfrac{4}{3}\right)} = -\dfrac{1}{7}$或$7$

        又因为直线$l$过点 $P\left( {1,2}\right)$，所以直线 $l$ 的方程为$7x-y-5=0$或$x+7y-15=0$
    \end{jieda}

    \item 若直线 ${2x} + y - 3 = 0$ 与直线 ${4x} + {2y} + a = 0$ 之间的距离不大于 $\sqrt{5}$ ,则 $a$ 的取值范围为\blkda{$[-16, 4]$}.
    \begin{jieda}
        $\dfrac{\left|\dfrac{a}{2}+3\right|}{\sqrt{5}} \leq \sqrt{5}$，解得$a \in [-16, 4]$
    \end{jieda}

    \item 下列各组点在直线 ${5x} - {12y} + 6 = 0$ 的异侧的是\dotfill{(\kh{C})}
    \xx{$(0,0)$与$(10,3)$}{$(3,5)$与$(6,7)$}{$(1,1)$与$(15,2)$}{$(-3,2)$与$(5,7)$}
    
     \item 设点 $P\left( {x,y}\right)$ 满足 ${ax} + {by} + c = 0$ ,则 “ $b = {2a}$ ” 是 “ $\left| {x + {2y} + 2}\right|  + \left| {x + {2y} - 1}\right|$ 为定值” 的\dotfill{(\kh{B})}
     \xx{充分不必要条件}{必要不充分条件}{充要条件}{既不充分也不必要条件}
     \begin{jieda}
         若 $\left| {x + {2y} + 2}\right|  + \left| {x + {2y} - 1}\right|  = \sqrt{5}\left( {\dfrac{\left| x + 2y + 2\right| }{\sqrt{5}} + \dfrac{\left| x + 2y - 1\right| }{\sqrt{5}}}\right)$ 为定值，即点 $P\left( {x,y}\right)$ 到直线 $x + {2y} + 2 = 0,x + {2y} - 1 = 0$ 两条直线距离之和为定值
        
        显然, 这两条直线平行, 如图,
        
        \begin{center}
        \includegraphics[width=0.3\textwidth]{bodyxb/src/bo_d2ohv9f7aajc738rr47g_12_241_207_363_276_0.jpg}
        \end{center}
        
        所以当点 $P\left( {x,y}\right)$ 在与这两条直线平行的直线上时,此时直线 ${ax} + {by} + c = 0$ 满足 ${ab} \neq  0$ 且 $b = {2a},$
        
        即 $b = {2a}$ ,且 $a \neq  0,b \neq  0,\left| {x + {2y} + 2}\right|  + \left| {x + {2y} - 1}\right|$ 为定值,
        
        所以 “ $b = {2a}$ ” 是 “ $\left| {x + {2y} + 2}\right|  + \left| {x + {2y} - 1}\right|$ 为定值” 的必要不充分条件.
     \end{jieda}

    \item 已知 $\bigtriangleup  {{ABC}}$ 三边所在直线分别为 ${AB} : x + {2y} - 4 = 0,{BC} : x - {2y} = 0,{CA} : x - y + 2 = 0$ .
    \begin{enumerate}
        \item 求 $\angle {BAC}$ 的角平分线所在直线的方程
        \item 求 $\bigtriangleup  {{ABC}}$ 的内心的坐标
        \item 求 $\bigtriangleup  {{ABC}}$ 的面积
    \end{enumerate}
    \begin{jieda}
        $A(0, 2), B(2, 1), C(-4, -2)$
        \begin{enumerate}
            \item  \begin{enumerate}
                \item 方法1

                设$\angle BAC$的角分线上点为$P(x, y)$，则有$\dfrac{|x + {2y} - 4|}{\sqrt{5}} = \dfrac{|x - y + 2|}{\sqrt{2}}$，故直线为$\dfrac{x + {2y} - 4}{\sqrt{5}} = \dfrac{x - y + 2}{\sqrt{2}}$或$\dfrac{x + {2y} - 4}{\sqrt{5}} = -\dfrac{x - y + 2}{\sqrt{2}}$

                数形结合可知$\angle {BAC}$ 的角平分线所在直线的方程的斜率小于0，因此直线方程为$\dfrac{x + {2y} - 4}{\sqrt{5}} = -\dfrac{x - y + 2}{\sqrt{2}}$，整理得$y =  - \left( {3 + \sqrt{10}}\right) x + 2$.

                \item 方法2

                $\vv{AB} = (2, -1)$，$\vv{AC} = (-4, -4)$，两者对应的单位向量为$\left(\dfrac{2\sqrt{5}}{5}, -\dfrac{\sqrt{5}}{5}\right)$和$\left(-\dfrac{\sqrt{2}}{2}, -\dfrac{\sqrt{2}}{2}\right)$，故角分线所在直线的方向向量为$\left(4\sqrt{5}-5\sqrt{2}, -(2\sqrt{5}+5\sqrt{2})\right)$，可得法向量为$(2\sqrt{5}+5\sqrt{2}, 4\sqrt{5}-5\sqrt{2})$，因此直线方程为$(2\sqrt{5}+5\sqrt{2})x+(4\sqrt{5}-5\sqrt{2})(y-2) = 0$，即$y =  - \left( {3 + \sqrt{10}}\right) x + 2$
            \end{enumerate}
            \item 容易判断$\angle ABC$的角分线为$y = 1$，故根据角分线的交点是内心可得内心的坐标为$\left( {\sqrt{10} - 3,1}\right)$
            \item 根据$\vv{AB} = (2, -1)$，$\vv{AC} = (-4, -4)$可知$S = \dfrac{1}{2}\left|2 \times (-4) - (-1)\times (-4)\right| = 6$.
        \end{enumerate}
    \end{jieda}
    
    \item 已知两条互相平行的直线分别过点 $M\left( {6,2}\right)$ 和 $N\left( {-3, - 1}\right)$ ,并且各自绕着点 ${M}$ 、 ${N}$ 以相同速度旋转. 在旋转过程中,求这两条平行线间的距离 $d$ 的取值范围,并求出距离 $d$ 最大时的直线方程.
    \begin{jieda}
        $d \in  (0,3\sqrt{10}\rbrack$
        
        过 ${M}$ 点直线方程为${3x} + y - {20} = 0$ ,过 ${N}$ 点直线方程为${3x} + y + {10} = 0$
    \end{jieda}
    
    \item 已知直线 $l$ 经过点 $P\left( {3,1}\right)$ ,且被两平行线 $\left\{  \begin{array}{l} {l}_{1} : x + y + 1 = 0 \\  {l}_{2} : {2x} + {my} + {12} = 0\left( {m \in  \mathbf{R}}\right)  \end{array}\right.$ 所截得的线段长为5. 
    \begin{enumerate}
        \item 求 ${l}_{1}$ , ${l}_{2}$ 之间的距离
        \item 求直线 $l$ 的方程
    \end{enumerate}
    \begin{jieda}
        \begin{enumerate}
            \item 当 ${l}_{1},{l}_{2}$ 平行时,则 $- 1 = \dfrac{-2}{m}$ ,解得 $m = 2$ ,此时 ${l}_{2} : x + y + 6 = 0$，则 ${l}_{1},{l}_{2}$ 之间的距离 $d = \dfrac{\left| 6 - 1\right| }{\sqrt{{1}^{2} + {1}^{2}}} = \dfrac{5\sqrt{2}}{2}$ 
            \item 设直线 $l$ 与直线 ${l}_{1},{l}_{2}$ 分别交于点 $A\left( {{x}_{1},{y}_{1}}\right) ,B\left( {{x}_{2},{y}_{2}}\right)$ ,
            
            则 $\left\{  \begin{array}{l} {x}_{1} + {y}_{1} + 1 = 0 \\  {x}_{2} + {y}_{2} + 6 = 0 \end{array}\right.$ ,两式相减得: $\left( {{x}_{1} - {x}_{2}}\right)  + \left( {{y}_{1} - {y}_{2}}\right)  = 5$ ,而 $\left| {AB}\right|  = 5$ ,
            
            即 ${\left( {x}_{1} - {x}_{2}\right) }^{2} + {\left( {y}_{1} - {y}_{2}\right) }^{2} = {25}$ ,解得 ${x}_{1} - {x}_{2} = 0$ 或 ${y}_{1} - {y}_{2} = 0$ ,
            
            由 ${x}_{1} - {x}_{2} = 0$ ,即 ${x}_{1} = {x}_{2},l \bot  x$ 轴,得直线 $l$ 方程为 $x = 3$ ,经验证符合题意,
            
            由 ${y}_{1} - {y}_{2} = 0$ ,即 ${y}_{1} = {y}_{2},l \bot  y$ 轴,得直线 $l$ 方程为 $y = 1$ ,经验证符合题意,
            
            所以直线 $l$ 的方程为 $x = 3$ 或 $y = 1$ .
           
            \begin{center}
            \includegraphics[width=0.3\textwidth]{bodyxb/src/bo_d2ohv9f7aajc738rr47g_13_240_498_370_352_0.jpg}
            \end{center}
        \end{enumerate}
    \end{jieda}

    
    \item 最近国际局势波云诡谲,我国在某地区进行军事演练,如图, $O,A,B$ 是三个军事基地, $C$ 为一个军事要塞,在线段 ${AB}$ 上. 已知 $\tan \angle {AOB} =  - 3,{OA} = {10}{{km}}$ , $C$ 到 ${OA}$ , ${OB}$ 的距离分别为 $5{{km}}$ , $2\sqrt{10}{{km}}$ . 以点 $O$ 为坐标原点,直线 ${OA}$ 为 $x$ 轴,建立平面直角坐标系如图所示, $C$ 位于第一象限. 
    
    \begin{center}
    \includegraphics[width=0.3\textwidth]{bodyxb/src/bo_d2ohv1v7aajc738rr44g_2_238_1155_377_337_0.jpg}
    \end{center}
    \begin{enumerate}
        \item 求两个军事基地 ${AB}$ 的长
        \item 若要塞 $C$ 正北方向距离要塞 ${10}{{km}}$ 处有一 $E$ 处正在进行爆破试验,爆炸波生成 $t{h}$ 时的半径为 $r = 5\sqrt{at}$ (参数 $a$ 为大于零的常数),爆炸波开始生成时,一飞行器以 ${30}\sqrt{2}{{km}}/{h}$ 的速度自基地 $A$ 开往基地 $B$ ,问参数 $a$ 控制在什么范围内时,爆炸波不会波及到飞行器的飞行.
    \end{enumerate}
    \begin{jieda}
        \begin{enumerate}
            \item 由题设得: $A\left( {{10},0}\right)$ ,直线 ${OB}$ 的方程为 ${3x} + y = 0$ , $C\left( {{x}_{0},5}\right) \left( {{x}_{0} > 0}\right)$ ,
            
            由 $\dfrac{\left| 3{x}_{0} + 5\right| }{\sqrt{{3}^{2} + {1}^{2}}} = 2\sqrt{10}$ ,且 ${x}_{0} > 0$ ,解得 ${x}_{0} = 5$ ,所以 $C\left( {5,5}\right)$ .
            
            所以直线 ${AC}$ 的方程为 $y =  - \left( {x - {10}}\right)$ ,即 $x + y - {10} = 0$ ,
            
            联立方程 $\left\{  \begin{array}{l} y =  - {3x} \\  x + y - {10} = 0 \end{array}\right.$ ,解得 $\left\{  \begin{array}{l} x =  - 5 \\  y = {15} \end{array}\right.$ ,即 $B\left( {-5,{15}}\right)$ ,
            
            所以 $\left| {AB}\right|  = \sqrt{{\left( -5 - {10}\right) }^{2} + {\left( {15} - 0\right) }^{2}} = {15}\sqrt{2}$，即基地 ${AB}$ 的长为 ${15}\sqrt{2}{{km}}$ .
            
            
            \item 设爆炸产生的爆炸波圆 $E$，由题意可得 $E\left( {5,{15}}\right)$ ,爆炸波生成 $t$ 小时后,飞行在线段 ${AB}$ 上的点 $F$ 处，则 $\left| {AF}\right|  = {30}\sqrt{2}t,0 \leq  t \leq  \dfrac{1}{2}$ ,所以 $F\left( {{10} - {30t},{30t}}\right)$
            
            爆炸波不会波及飞行器的通行,即 $\left| {EF}\right|  > r$ 对 $t \in  \left\lbrack  {0,\dfrac{1}{2}}\right\rbrack$ 恒成立.
            
            所以 ${\left| EF\right| }^{2} = {\left( {30}t - 5\right) }^{2} + {\left( {30}t - {15}\right) }^{2} > {r}^{2} = {25at}$ ,即 ${\left( {30}t - 5\right) }^{2} + {\left( {30}t - {15}\right) }^{2} > {25at}$

            \begin{dingautolist}{172}
                \item 当 $t = 0$ 时,上式恒成立
                \item 当 $t \in  \left( {0,\dfrac{1}{2}}\right\rbrack$ 时,整理得 $a < {72t} + \dfrac{10}{t} - {48}$ 
                
                因为 ${72t} + \dfrac{10}{t} - {48} \geq  2\sqrt{{72t} \times  \dfrac{10}{t}} - {48} = {24}\sqrt{5} - {48}$，当且仅当 ${72t} = \dfrac{10}{t}$ ,即 $t = \dfrac{\sqrt{5}}{6}$ 时,等号成立
            \end{dingautolist}

            所以在 $0 < a < {24}\sqrt{5} - {48}$ 时, $r < \left| {EF}\right|$ 恒成立,亦即爆炸波不会波及飞行的通行.
        \end{enumerate}
    \end{jieda}

    \item 已知直线 $l : {ax} + {by} + c = 0$ 和点 $P\left( {{x}_{0},{y}_{0}}\right)$ ,点 $P$ 到直线 $l$ 的有向距离 $d\left( {P,l}\right)$ 用如下方法规定: 若 $b \neq  0,d\left( {P,l}\right)  = \dfrac{\left| b\right| \left| {a{x}_{0} + b{y}_{0} + c}\right| }{b\sqrt{{a}^{2} + {b}^{2}}}$ ,若 $b = 0,d\left( {P,l}\right)  = \dfrac{a{x}_{0} + c}{a}$ .
    \begin{enumerate}
        \item 已知直线 ${l}_{1} : {3x} - {4y} + {12} = 0$ ,直线 ${l}_{2} : {2x} + 3 = 0$ ,求原点 $O$ 到直线 ${l}_{1},{l}_{2}$ 的有向距离 $d\left( {O,{l}_{1}}\right) ,d\left( {O,{l}_{2}}\right)$
        \item 已知点 $A\left( {2,1}\right)$ 和点 $B\left( {3, - 1}\right)$ ,是否存在通过点 ${A}$ 的直线 ${l}_{3}$ ,使得 $d\left( {B,{l}_{3}}\right)  = 2$ ? 如果存在,求出所有这样的直线 ${l}_{3}$ ,如果不存在,说明理由
        \item 设直线 ${l}_{4} : x\cos a + {2y}\sin a - 2 = 0$ ,问是否存在实数 $t > 0$ ,使得对任意的参数 $a$ 都有: 点 ${F}_{1}\left( {-t,0}\right) ,{F}_{2}\left( {t,0}\right)$ 到 ${l}_{4}$ 的有向距离 $d\left( {{F}_{1},{l}_{4}}\right) ,d\left( {{F}_{2},{l}_{4}}\right)$ 满足 $d\left( {{F}_{1},{l}_{4}}\right)  \cdot  d\left( {{F}_{2},{l}_{4}}\right)  = 1$ ? 如果满足,求出所有满足条件的实数 $t$ ; 如果不存在,请说明理由.
    \end{enumerate}
    \begin{jieda}
        \begin{enumerate}
            \item 由直线 ${l}_{1} : {3x} - {4y} + {12} = 0$ ,直线 ${l}_{2} : {2x} + 3 = 0$ ,根据点到直线的有向距离公式得, $d\left( {O,{l}_{1}}\right)  = \dfrac{\left| {-4}\right| \left| {0 + 0 + {12}}\right| }{-4\sqrt{{3}^{2} + {\left( -4\right) }^{2}}} =  - \dfrac{12}{5},d\left( {O,{l}_{2}}\right)  = \dfrac{2 \cdot  0 + 3}{2} = \dfrac{3}{2}$，即 $d\left( {O,{l}_{1}}\right)  =  - \dfrac{12}{5},d\left( {O,{l}_{2}}\right)  = \dfrac{3}{2}$
            
            \item 
            \begin{dingautolist}{172}
                \item 当直线 ${l}_{3}$ 的斜率不存在时,直线 ${l}_{3}$ 的方程为 $x - 2 = 0$，此时 $d\left( {B,{l}_{3}}\right)  = \dfrac{1 \times  3 - 2}{1} = 1 \neq  2$ ,舍去
                \item 当直线 ${l}_{3}$ 的斜率存在时,直线 ${l}_{3}$ 的方程为 $y - 1 = k\left( {x - 2}\right)$ ,
                
                由题意 $b > 0$ ,所以直线可化为 $- {kx} + y - 1 + {2k} = 0$ ,
                
                假设 $d\left( {B,{l}_{3}}\right)  = \dfrac{\left| 1\right| \left| {-{3k} - 1 - 1 + {2k}}\right| }{1 \times  \sqrt{{k}^{2} + 1}} = 2$ ,则 $3{k}^{2} - {4k} = 0$ ,解得 $k = 0$ 或 $\dfrac{4}{3}$ .
                \end{dingautolist}
            所以存在直线 ${l}_{3}$ 的方程为 $y - 1 = 0$ 或 ${4x} - {3y} - 5 = 0$ 
            
            \item 
            \begin{dingautolist}{172}
                \item 当 $\sin a = 0$ 时,直线 ${l}_{4} : x\cos a - 2 = 0$，$d\left( {{F}_{1},{l}_{4}}\right)  = \dfrac{-t\cos a - 2}{\cos a},d\left( {{F}_{2},{l}_{4}}\right)  = \dfrac{t\cos a - 2}{\cos a}$ ,
                
                由 $d\left( {{F}_{1},{l}_{4}}\right)  \cdot  d\left( {{F}_{2},{l}_{4}}\right)  = \dfrac{\left( {-t\cos a - 2}\right) \left( {t\cos a - 2}\right) }{{\cos }^{2}a} = 1$ ,整理得$4 - {t}^{2}{\cos }^{2}a = {\cos }^{2}a$
                
                $\because {\sin }^{2}a + {\cos }^{2}a = 1,\therefore {t}^{2} = 3,\because t > 0$ ,即 $t = \sqrt{3}$
                \item 当 $\sin a \neq  0$ 时,直线 ${l}_{4} : x\cos a + {2y}\sin a - 2 = 0$ ,
            
                得 $d\left( {{F}_{1},{l}_{4}}\right)  = \dfrac{\left| {2\sin a}\right| \left| {-t\cos a - 2}\right| }{\left( {2\sin a}\right) \sqrt{{\cos }^{2}a + 4{\sin }^{2}a}},d\left( {{F}_{2},{l}_{4}}\right)  = \dfrac{\left| {2\sin a}\right| \left| {t\cos a - 2}\right| }{\left( {2\sin a}\right) \sqrt{{\cos }^{2}a + 4{\sin }^{2}a}}$ ,
                
                由 $d\left( {{F}_{1},{l}_{4}}\right)  \cdot  d\left( {{F}_{2},{l}_{4}}\right)  = \dfrac{\left| \left( t\cos a - 2\right) \left( -t\cos a - 2\right) \right| }{{\cos }^{2}a + 4{\sin }^{2}a} = 1$ ,
                
                即 $\left| {4 - {t}^{2}{\cos }^{2}a}\right|  = {\cos }^{2}a + 4{\sin }^{2}a = 4 - 3{\cos }^{2}a$，$4 - {t}^{2}{\cos }^{2}a = 4 - 3{\cos }^{2}a$ 或 $4 - {t}^{2}{\cos }^{2}a = 3{\cos }^{2}a - 4$ ,解得${t}^{2} = 3$ 或 $\left( {{t}^{2} + 3}\right) {\cos }^{2}a = 8$
                
                由题意对任意的参数 $a$ 都有 $d\left( {{F}_{1},{l}_{4}}\right)  \cdot  d\left( {{F}_{2},{l}_{4}}\right)  = 1$ 恒成立,所以 $t = \sqrt{3}$
            \end{dingautolist}
            综上所述,存在实数 $t > 0$ 满足题目条件,即 $t = \sqrt{3}$
        \end{enumerate}
    \end{jieda}

\end{enumerate}

C 组：
\begin{enumerate}
    \item 平面中两条直线 ${l}_{1}$ 、 ${l}_{2}$ 相交于点 $O$ ,对于平面上任意一点 $M$ ,若 $p$ , $q$ 分别是 $M$ 到直线 ${l}_{1}$ 和 ${l}_{2}$ 的距离,则称有序非负实数对(p, q)是点 $M$ 的“距离坐标”. 已知常数 $p \geq  0,q \geq  0$ , 给出下列命题:
    \begin{dingautolist}{172}
        \item 若 $p = q = 0$ ,则 “距离坐标” 为$(0,0)$的点有且仅有 1 个
        \item 若 ${pq} = 0,p + q \neq  0$ ,则 “距离坐标” 为$(p, q)$的点有且仅有 2 个
        \item 若 ${pq} \neq  0$ ,则 “距离坐标” 为$(p, q)$的点有且仅有 4 个
    \end{dingautolist}
    以上命题中,正确的命题是\blkda{\ding{172}\ding{173}\ding{174}}.
    \begin{jieda}
        \begin{dingautolist}{172}
            \item 只有当点 $M$ 与点 $O$ 重合时,满足题意. 故正确
            
            \item 不妨假设 $p = 0$ ,则点 $M$ 在直线 ${l}_{1}$ 上. 又点 $M$ 到直线 ${l}_{2}$ 的距离为 $q \neq  0$ ,

        \begin{center}
        \includegraphics[width=0.5\textwidth]{bodyxb/src/bo_d2ohv9f7aajc738rr47g_17_246_196_599_479_0.jpg}
        \end{center}
        
        如图，作 ${l}_{2}$ 的两条平行线 ${l}_{21},{l}_{22}$ ,使得 ${l}_{21},{l}_{22}$ 与 ${l}_{2}$ 的距离均为 $q$ .
        
        由定义结合图象可知,直线 ${l}_{21},{l}_{22}$ 与 ${l}_{1}$ 的交点满足条件.
        
        所以,“距离坐标”为$(p, q)$的点有且仅有 2 个. 故正确;
        
        \item 如图，作 ${l}_{1}$ 的两条平行线 ${l}_{11},{l}_{12}$ ,使得 ${l}_{11},{l}_{12}$ 与 ${l}_{1}$ 的距离均为 $p$ ；作 ${l}_{2}$ 的两条平行线 ${l}_{21},{l}_{22}$ ,使得 ${l}_{21},{l}_{22}$ 与 ${l}_{2}$ 的距离均为 $q$ .
        
        \begin{center}
        \includegraphics[width=0.5\textwidth]{bodyxb/src/bo_d2ohv9f7aajc738rr47g_17_249_1111_627_478_0.jpg}
        \end{center}
        
        由定义结合图象可知, 满足条件的点有且仅有 4 个. 故正确.  
        \end{dingautolist}
    \end{jieda}
    
    \item 若恰有三组不全为 0 的实数对$(a, b)$满足关系式 $\left| {{2a} + b + 1}\right|  = \left| {a - {5b} - 1}\right|  = t\sqrt{{a}^{2} + {b}^{2}}$ ,写出符合条件实数 $t$ 的取值集合\blkda[8]{$\dfrac{5}{2}$ 或 $\dfrac{{11}\sqrt{37}}{37}$ 或 $\dfrac{11}{5}$ (写出其中一个即可)}.
    \begin{jieda}
        由 $\left| {{2a} + b + 1}\right|  = \left| {a - {5b} - 1}\right|  = t\sqrt{{a}^{2} + {b}^{2}}$，得 $\dfrac{\left| 2a + b + 1\right| }{\sqrt{{a}^{2} + {b}^{2}}} = \dfrac{\left| -a + 5b + 1\right| }{\sqrt{{a}^{2} + {b}^{2}}} = t$，则 ${d}_{1} = \dfrac{\left| 2a + b + 1\right| }{\sqrt{{a}^{2} + {b}^{2}}}$ 表示点 $A\left( {2,1}\right)$ 到直线 ${ax} + {by} + 1 = 0$ 的距离；
        
        ${d}_{2} = \dfrac{\left| -a + 5b + 1\right| }{\sqrt{{a}^{2} + {b}^{2}}}$ 表示点 $B\left( {-1,5}\right)$ 到直线 ${ax} + {by} + 1 = 0$ 的距离；
        
        直线 ${ax} + {by} + 1 = 0$ 不过原点, $\left| {AB}\right|  = 5$

        \begin{dingautolist}{172}
            \item 当 $t = \dfrac{\left| AB\right| }{2} = \dfrac{5}{2}$ ,此时,易得符合题意的直线 $l$ 为线段 ${AB}$ 的垂直平分线 ${6x} - {8y} + {21} = 0$ 以及直线 ${AB}$ 平行的两条直线 ${8x} + {6y} + 3 = 0$ 和 ${8x} + {6y} - {47} = 0$ 
            
            \item 当 $t < \dfrac{\left| AB\right| }{2} = \dfrac{5}{2}$ 时,有 4 条直线 $l$ 会使得点 $A\left( {2,1}\right)$ 和 $B\left( {-1,5}\right)$ 到它们的距离相等,
            
            注意到 $l$ 不过原点,所以,当其中一条直线过原点时,会作为增根被舍去.
            
            设点 $A$ 到 $l$ 的距离为 $d$ .

            \begin{enumerate}
                \item 作为增根被舍去的直线 $l$ ,过原点和 $A,B$ 的中点 $M\left( {\dfrac{1}{2},3}\right)$ ,其方程为 ${6x} - y = 0$ ,此时, $t = d = \dfrac{{11}\sqrt{37}}{37} < \dfrac{5}{2}$ ,符合
                \item 作为增根被舍去的直线 $l$ ,过原点且与 ${AB}$ 平行,其方程为 ${4x} + {3y} = 0$，此时, $t = d = \dfrac{11}{5} < \dfrac{5}{2}$ ,符合.
            \end{enumerate}
            \item 当 $t > \dfrac{\left| AB\right| }{2} = \dfrac{5}{2}$ 时,至多只有两条直线符合题意,舍去.
        
        \end{dingautolist}
        综上, $t = \dfrac{5}{2}$ 或 $\dfrac{{11}\sqrt{37}}{37}$ 或 $\dfrac{11}{5}$ .
    \end{jieda}
\end{enumerate}
\clearpage
\section*{直线专题: 直线中的定点和对称}

A组:

\begin{enumerate}
    \item 点$(-1,2)$关于直线 ${3x} - {5y} - {21} = 0$ 的对称点的坐标是\blkda{$(5, -8)$}.
    \begin{jieda}
        设对称点坐标为$(a, b)$，则有$\left\{\begin{array}{l}
             3 \times \dfrac{-1+a}{2} - 5 \times \dfrac{b+2}{2} - {21} = 0 \\
             \dfrac{a+1}{3} = \dfrac{b-2}{-5}
        \end{array}\right.$，解得$\left\{\begin{array}{l}
             a = 5  \\
             b = -8
        \end{array}\right.$
    \end{jieda}
    
    \item 已知点 $P\left( {3,2}\right)$ 与点 $Q\left( {1,4}\right)$ 关于直线 $l$ 对称,则直线 $l$ 的方程为\blkda[3]{$x - y + 1 = 0$}.
    \begin{jieda}
        设直线 $l$ 的斜率为 $k$ ,则 $k \times  \dfrac{4 - 2}{1 - 3} =  - 1 \Longrightarrow  k = 1$，直线 ${PQ}$ 的中点坐标为 $\left\{  \begin{array}{l} {x}_{0} = \dfrac{3 + 1}{2} = 2 \\  {y}_{0} = \dfrac{2 + 4}{2} = 3 \end{array}\right.$，所以由点斜式写出直线 $l$ 方程为 $y - 3 = x - 2$ ,即 $x - y + 1 = 0$ .
    \end{jieda}
    
    \item 已知直线 ${l}_{1}$ 与直线 ${l}_{2} : {2x} + {3y} - 6 = 0$ 关于点 $A\left( {1, - 1}\right)$ 对称,则直线 ${l}_{1}$ 的方程是\blkda[3]{$2x+3y+8=0$}.
    \begin{jieda}
        在直线$l_2$上取点$B(0,2)$，$C(3, 0)$，得到关于$A$的对称点$B'(2, -4)$，$C'(-1, -2)$，故得到直线$l_1$的方程：$\dfrac{x-(-1)}{3} = \dfrac{y-(-2)}{-2}$，即$2x+3y+8=0$.
    \end{jieda}
    
    \item 已知直线 ${l}_{1}$ 与直线 ${l}_{2} : {2x} - y - 2 = 0$ 关于直线 ${l}_{3} : x - y - 1 = 0$ 对称,则直线 ${l}_{1}$ 的方程是\blkda[3]{$x-2y-1=0$}.
    \begin{jieda}
        $l_2$与$l_3$交于$A(1, 0)$，故$l_1$过点$A$，在直线$l_2$上取点$B(0, -2)$，可求得其关于$l_3$的对称点$C(-1, -1)$，故$l_1$即直线$AC$：$\dfrac{x+1}{2} = \dfrac{y+1}{1}$，即$x-2y-1=0$.
    \end{jieda}
    
    \item 不论 $m$ 取何实数,直线 $\left( {{2m} - 1}\right) x - \left( {m + 3}\right) y - m + {11} = 0$ 恒过一个定点,则这个定点坐标是\blkda{$(2, 3)$}.
    \begin{jieda}
        $(2x-y-1)m+(11-x-3y) = 0$，$\left\{\begin{array}{l}
             2x-y-1 = 0 \\
             11-x-3y = 0
        \end{array}\right.$，解得$\left\{\begin{array}{l}
             x = 2 \\
             y = 3
        \end{array}\right.$
    \end{jieda}
    
    \item 点 $P\left( {{x}_{0},{y}_{0}}\right)$ 关于直线 $l : x + y + c = 0$ 的对称点的坐标为\blkda[3]{$(-y_0-c, -x_0-c)$}.
    
    \item 点 $P\left( {{x}_{0},{y}_{0}}\right)$ 关于直线 $l : x - y + c = 0$ 的对称点的坐标为\blkda[3]{$(y_0-c, x_0+c)$}.
\end{enumerate}

B组：
\begin{enumerate}
    \item 直线 ${3x} - {2y} = 0$ 关于点 $\left( {\dfrac{1}{3},0}\right)$ 对称的直线方程\blkda[3]{${3x} - {2y} - 2 = 0$}.
    \begin{jieda}
        设所求直线上任一点为$(x,y)$,则其关于点 $\left( {\dfrac{1}{3},0}\right)$ 对称的点为 $\left( {\dfrac{2}{3} - x, - y}\right)$，因为点 $\left( {\dfrac{2}{3} - x, - y}\right)$ 在直线 ${3x} - {2y} = 0$ 上，所以 $3\left( {\dfrac{2}{3} - x}\right)  - 2\left( {-y}\right)  = 0$ ,化简得 ${3x} - {2y} - 2 = 0$，所以所求直线方程为 ${3x} - {2y} - 2 = 0$
    \end{jieda}

    \item 点 $M\left( {-3,4}\right)$ 与点 $N\left( {1,0}\right)$ 关于直线 $l : {ax} - {by} + 3 = 0$ 对称,则 $a + b$ 的值为\blkda{2}. 
    \begin{jieda}
        因为 ${k}_{MN} = \dfrac{4 - 0}{-3 - 1} =  - 1$ ,故 $\dfrac{a}{b} = 1$ ,而 ${MN}$ 的中点为$(-1,2)$，故 $- a - {2b} + 3 = 0$ ,所以 $a = b = 1$ ,所以 $a + b = 2$
    \end{jieda}
    
    \item 若 $M\left( {m,n}\right)$ 、 $N\left( {n + 1,m - 1}\right)$ 两点关于直线 $l$ 对称,则直线 $l$ 的方程是\blkda[3]{$x-y-1=0$}.
    \begin{jieda}
        $MN$中点坐标$\left(\dfrac{m+n+1}{2}, \dfrac{m+n-1}{2}\right)$，$k_{MN} = \dfrac{n-m+1}{m-n-1}$，故$k_l = \dfrac{n-m+1}{n-m+1} = 1$，由此可得直线$l$的方程为$x-y-1=0$.
    \end{jieda}

    \item 已知 ${4A} + {5B} + 6 = 0$ . 若直线 $l : {Ax} + {By} + 1 = 0$ 必过定点,则这个定点的坐标是\blkda{$\left(\dfrac{2}{3}, \dfrac{5}{6}\right)$}.

    \item 已知直线 ${l}_{1} : {2x} + y + 2 = 0$ 与 ${l}_{2} : {4x} + {by} + c = 0$ 关于点 $P\left( {1,0}\right)$ 对称,则$b + c =$ \blkda{$-10$}.
    \begin{jieda}
        在直线 ${l}_{1} : {2x} + y + 2 = 0$ 上取点 $M\left( {-1,0}\right) ,N\left( {0, - 2}\right) ,M,N$ 关于点 $P\left( {1,0}\right)$ 对称的点分别为 ${M}_{1}\left( {3,0}\right) ,{N}_{1}\left( {2,2}\right)$ .
        
        点 ${M}_{1}\left( {3,0}\right) ,{N}_{1}\left( {2,2}\right)$ 在直线 ${l}_{2} : {4x} + {by} + c = 0$ 上,
        
        $\therefore {12} + c = 0,8 + {2b} + c = 0$ ,解得 $c =  - {12},b = 2$，$\therefore b + c =  - {10}$ .
    \end{jieda}

    \item 在平面直角坐标系内,一束光线从点 $A\left( {1,2}\right)$ 出发,被直线 $y = x$ 反射后到达点 $B\left( {3,6}\right)$ ,则这束光线从 $A$ 到 $B$ 所经过的距离为\blkda{$\sqrt{26}$}.
    \begin{jieda}
        作出点 $A$ 关于直线 $y = x$ 的对称点 $C\left( {2,1}\right)$

        \begin{center}
            \includegraphics[width=0.2\linewidth]{bodyxb//src/zhixianduichen1.png}
        \end{center}
        
        连接 ${CB}$ ,交直线 $y = x$ 于点 $M$ ,则 $\left| {CB}\right|$ 即为光线经过路程的最小值,且$\left| {CB}\right|  = \sqrt{{\left( 3 - 2\right) }^{2} + {\left( 6 - 1\right) }^{2}} = \sqrt{26}$ ,此即光线从 $A$ 到 $B$ 所经过的距离为 $\sqrt{26}$ .
    \end{jieda}

    \item 费马点是指三角形内到三角形三个顶点距离之和最小的点, 当三角形三个内角均小$\ang{120}$时,费马点与三个顶点连线正好三等分费马点所在的周角,即该点所对三角形三边的张角相等,均为 ${120}^{ \circ  }$ . 根据以上性质,已知 $A\left( {-2,0}\right)$ , $B\left( {2,0}\right)$ ,$C\left( {0,4}\right)$ , $P$ 为 $\bigtriangleup  {{ABC}}$ 内一点,记 $f\left( P\right)  = \left| {PA}\right|  + \left| {PB}\right|  + \left| {PC}\right|$ ,则 $f\left( P\right)$ 的最小值为\blkda{$4 + 2\sqrt{3}$}.
    \begin{jieda}
        设 $O\left( {0,0}\right)$ 为坐标原点,由 $A\left( {-2,0}\right) ,B\left( {2,0}\right) ,C\left( {0,4}\right)$，知 $\left| {AC}\right|  = \left| {BC}\right|  = 2\sqrt{5}$ ,且 $\vee  {ABC}$ 为锐角三角形，因此,费马点 $M$ 在线段 ${OC}$ 上,设 $M\left( {0,h}\right)$ 
        \begin{center}
            \includegraphics[width=0.2\linewidth]{bodyxb//src/zhixianduichen2.png}
        \end{center}
        则 $\bigtriangleup {MAB}$ 为顶角是 ${120}^{ \circ  }$ 的等腰三角形,故 $h = \left| {OB}\right| \tan {30}^{ \circ  } = \dfrac{2\sqrt{3}}{3}$，所以 $f\left( P\right)  \geq  f\left( M\right)  = \left| {MA}\right|  + \left| {MB}\right|  + \left| {MC}\right|  = {4h} + 4 - h = 4 + 2\sqrt{3}$，则 $f\left( P\right)$ 的最小值为 $4 + 2\sqrt{3}$ .
    \end{jieda}

    \item 将一张长为 2 宽为 1 的矩形纸片 ${OABC}$ 置于平面直角坐标系中,其中$O\left( {0,0}\right) ,A\left( {2,0}\right) ,C\left( {0,1}\right)$ ,将矩形纸片折叠,使 $O$ 点落在线段 ${BC}$ 上,设折痕所在直线的斜率为 $k$ ,则 $k$ 的取值范围是\blkda{$\left\lbrack  {-2,0}\right\rbrack$}.
    \begin{jieda}
        要想使折叠后 $O$ 点落在线段 ${BC}$ 上,可取 ${BC}$ 上任意一点 $D$ ,作线段 ${OD}$ 的垂直平分线 $l$ ,以$l$ 为折痕可使 $O$ 与 $D$ 重合,
        因为 ${k}_{OD} \geq  {k}_{OB} = \dfrac{1}{2}$ ,所以 $k =  - \dfrac{1}{{k}_{OD}} \geq   - 2$ ,且 $k < 0$ .
        又当折叠后 $O$ 与 $C$ 重合时, $k = 0$ ,所以 $- 2 \leq  k \leq  0,\therefore k$ 的取值范围是 $\left\lbrack  {-2,0}\right\rbrack$ .
    \end{jieda}
    
    \item 已知点 $M\left( {a,b}\right)$ 与点 $N\left( {0, - 1}\right)$ 在直线 ${3x} - {4y} + 5 = 0$ 的两侧,给出以下结论:
    \begin{dingautolist}{172}
        \item ${3a} - {4b} + 5 > 0$
        \item 当 $a > 0$ 时, $a + b$ 有最小值,无最大值
        \item ${a}^{2} + {b}^{2} > 1$
        \item 当 $a > 0$ 且 $a \neq  1$ 时, $\dfrac{b + 1}{a - 1}$ 的取值范围是 $\left( {-\infty , - \dfrac{9}{4}}\right)  \cup  \left( {\dfrac{3}{4}, + \infty }\right)$
    \end{dingautolist}
    正确的个数是\dotfill{(\kh{B})}
    \xx{1}{2}{3}{4}
    \begin{dingautolist}{172}
        \item 将 $N\left( {0, - 1}\right)$ 代入有 $3 \times  0 - 4 \times  \left( {-1}\right)  + 5 = 9 > 0$，而 $M$ 与 $N$ 在 ${3x} - {4y} + 5 = 0$ 的两侧,则 ${3a} - {4b} + 5 < 0$，故错误;
        \item 由上知: ${3a} - {4b} + 5 < 0$ 且 $a > 0$ ,则根据图可知 $M$ 在直线上方与 $y$ 轴右侧部分，所以 $a + b > \dfrac{5}{4}$ ,故 $a + b$ 无最值，故错误
        \begin{center}
            \includegraphics[width=0.3\linewidth]{bodyxb//src/zhixianduichen3.png}
        \end{center}
        \item 由上图知: $M$ 在直线左上方,则 ${a}^{2} + {b}^{2} > {\left( \dfrac{5}{\sqrt{{3}^{2} + {4}^{2}}}\right) }^{2} = 1$，故正确;
        
        \item 由 ${3x} - {4y} + 5 = 0$ 过 $\left( {0,\dfrac{5}{4}}\right)$ 且 $a > 0$ 且 $a \neq  1$ ,即 $M$ 在直线上方与 $y$ 轴右侧部分，而 $\dfrac{b + 1}{a - 1}$ 表示$(1, - 1)$与 $M$ 连线的斜率,由图知: $\dfrac{b + 1}{a - 1} \in  \left( {-\infty , - \dfrac{9}{4}}\right)  \cup  \left( {\dfrac{3}{4}, + \infty }\right)$ ,故正确.
    \end{dingautolist}
    
    \item 已知画有平面直角坐标系的一张纸上有两点 $A\left( {-1, - 4}\right)$ 和 $B\left( {3,0}\right)$ ,若把这张纸折叠起来,使 $A,B$ 两点重合.
    \begin{enumerate}
        \item 求折痕所在的直线的方程
        \item 求折叠后与坐标原点重合的点的坐标
    \end{enumerate}
    \begin{jieda}
        \begin{enumerate}
            \item 即求$A,B$两点的对称轴的直线，$AB$中点为$(1, -2)$，$k_{AB} = 1$，故$k_l = -1$，因此折痕$l$的方程为：$x+y+1=0$
            \item $(-1, -1)$
        \end{enumerate}
    \end{jieda}
    
    \item 已知 $a,b$ 均为实数,且 ${2a} - {3b} = 1$ ,直线 $l : {ax} + {by} - 5 = 0$ 是否过定点? 若过定点, 求出这个点的坐标; 若不过定点, 请说明理由.
    \begin{jieda}
        ${2a} - {3b} = 1 \Longleftrightarrow {10a} - {15b} = 5$，因此直线过定点$(10, -15)$
    \end{jieda}
    
    \item 求 证: 对于直线 ${l}_{1} : \left( {{3k} + 1}\right) x + \left( {1 - {4k}}\right) y = {5k} + 4$ 和直 线 ${l}_{2} : \left( {k + 1}\right) x + \left( {1 - {2k}}\right) y = {3k}$ ,不论 $k$ 取何实数,都存在一条确定的直线与 ${l}_{1}$ 和 ${l}_{2}$ 都有公共点, 并求出这条直线的方程.
    \begin{jieda}
        $l_1$过定点$(3, 1)$，$l_2$过定点$(1, -1)$，因此目标直线应过这两个定点，方程为$-x+y+2=0$
    \end{jieda}

    \item 在 $\bigtriangleup {ABC}$ 中,已知顶点 $A$ 的坐标是 $\left( {3, - 1}\right) ,{AB}$ 边上的中线 ${CM}$ 所在直线方程为: ${6x} + {10y} - {59} = 0,\angle B$ 的平分线方程 ${BT}$ 为: $x - {4y} + {10} = 0$ ,求直线 ${BC}$ 的方程.
    \begin{jieda}
        设$M\left(a, \dfrac{59-6a}{10}\right)$，则$B\left(2a-3, \dfrac{59-6a}{5}+1\right)$，则有$(2a-3)-4\left(\dfrac{64-6a}{5}\right) + 10 = 0$，解得$a = \dfrac{13}{2}$，故$B\left(10, 5\right)$，$A$关于$x - {4y} + {10} = 0$的对称点为$A'(1, 7)$，故直线$BC$即直线$BA': 2x+9y-65=0$
    \end{jieda}
    
    
    \item \begin{enumerate}
        \item 在直线 ${3x} - y - 1 = 0$ 上求一点 $M$ ,使它到点 ${A}\left( {4,1}\right)$ 和 ${B}\left( {3,0}\right)$ 的距离之和最小, 并求最小值
        \item 在直线 ${3x} - y - 1 = 0$ 上求一点 $M$ ,使它到点 ${A}\left( {4,1}\right)$ 和 ${B}\left( {0,4}\right)$ 的距离之差最大,并求最大值
        \item 已知点 ${M}\left( {3,5}\right)$ ,在直线 $l : x - {2y} + 2 = 0$ 和 $y$ 轴上各找一点 ${P}$ 和 ${Q}$ ,使 $\bigtriangleup  {{MPQ}}$ 的周长最小
    \end{enumerate}
    \begin{jieda}
        \begin{enumerate}
            \item $M\left( {\dfrac{7}{9},\dfrac{4}{3}}\right) ,\left| {AM}\right|  + \left| {BM}\right|  = \sqrt{34}$
            
            \item $M\left( {2,5}\right) ,\parallel {AM}\left| -\right| {BM}\parallel  = \sqrt{5}$
            
            \item $P\left( {\dfrac{5}{2},\dfrac{9}{4}}\right) ,Q\left( {0,\dfrac{7}{2}}\right)$
        \end{enumerate}
    \end{jieda}
    
    \item 一束光从光源 $C\left( {1,2}\right)$ 射出,经 $x$ 轴反射后(反射点为 $M$ ),射到线段 $y =  - x + b,x \in  \left\lbrack  {3,5}\right\rbrack$ 上 $N$ 处.
    \begin{enumerate}
        \item 若 $M\left( {3,0}\right) ,b = 7$ ,求光从 $C$ 出发,到达点 $N$ 时所走过的路程
        \item 若 $b = 8$ ,求反射光的斜率的取值范围
        \item 若 $b \geq  6$ ,求光从 $C$ 出发,到达点 $N$ 时所走过的最短路程
    \end{enumerate}
    \begin{jieda}
        \begin{enumerate}
            \item $C\left( {1,2}\right)$ 关于 $x$ 轴的对称点 ${C}^{\prime }\left( {1, - 2}\right) ,{l}_{{C}^{\prime }M} : y = x - 3$，由 $\left\{  {\begin{array}{l} y = x - 3 \\  y =  - x + 7 \end{array} \Longrightarrow  x = 5 \in  \left\lbrack  {3,5}\right\rbrack  }\right.$ ,则此时 $N\left( {5,2}\right)$，所以光所走过的路程即 $\left| {{C}^{\prime }N}\right|  = \sqrt{{\left( 5 - 1\right) }^{2} + {\left( 2 + 2\right) }^{2}} = 4\sqrt{2}$ .
            \item 对于线段 $y =  - x + 8,x \in  \left\lbrack  {3,5}\right\rbrack$ ,令其端点 $A\left( {3,5}\right) ,B\left( {5,3}\right)$，则 ${k}_{{C}^{\prime }A} = \dfrac{7}{2},{k}_{{C}^{\prime }B} = \dfrac{5}{4}$，所以反射光斜率的取值范围是 $\left\lbrack  {\dfrac{5}{4},\dfrac{7}{2}}\right\rbrack$ .
            \item 若反射光与直线 $y =  - x + b$ 垂直,则由 $\left\{  {\begin{array}{l} y =  - x + b \\  y = x - 3 \end{array} \Longrightarrow  x = \dfrac{b + 3}{2}}\right.$ .
            \begin{dingautolist}{172}
                \item 当 $x = \dfrac{b + 3}{2} \in  \left\lbrack  {3,5}\right\rbrack$ ,即 $6 \leq  b \leq  7$ 时,光所走过的最短路程为点 ${C}^{\prime }$ 到直线 $y =  - x + b$ 的距离
                
                所以路程 $S = \dfrac{\left| 1 - 2 - b\right| }{\sqrt{2}} = \dfrac{b + 1}{\sqrt{2}}$ .

                \item 当 $x = \dfrac{b + 3}{2} \in  \lbrack 5, + \infty )$ ,即 $b > 7$ 时,光所走过的最短路程为线段 ${C}^{\prime }B$ ,其中 $B\left( {5,b - 5}\right)$，
                所以 $\left| {{C}^{\prime }B}\right|  = S = \sqrt{{\left( 5 - 1\right) }^{2} + {\left( b - 3\right) }^{2}} = \sqrt{{b}^{2} - {6b} + {25}}$ .
                综上: $S = \left\{  \begin{array}{l} \dfrac{b + 1}{\sqrt{2}},6 \leq  b \leq  7 \\  \sqrt{{b}^{2} - {6b} + {25}},b > 7 \end{array}\right.$ .
            \end{dingautolist}
        \end{enumerate}
    \end{jieda}
\end{enumerate}

C 组：
\begin{enumerate}
    \item 若恰有三组不全为 0 的实数对(a, b)满足关系式 $\left| {a + b + 1}\right|  = \left| {{4a} - {3b} + 1}\right|  = t\sqrt{{a}^{2} + {b}^{2}}$ , 则实数 $t$ 的所有可能的值为\blkda[3]{$\dfrac{5}{2}$ 或 $\dfrac{7}{5}$ 或 $\dfrac{7\sqrt{29}}{29}$}.
    \begin{jieda}
        由已知得 $t > 0$ ,整理得 $\dfrac{\left| a + b + 1\right| }{\sqrt{{a}^{2} + {b}^{2}}} = \dfrac{\left| 4a - 3b + 1\right| }{\sqrt{{a}^{2} + {b}^{2}}} = t$，看成有且仅有三条直线满足 $A\left( {1,1}\right)$ 和 $B\left( {4, - 3}\right)$ 到直线 $l : {ax} + {by} + 1 = 0$ (不过原点)的距离 $t$ 相等,又 $\left| {AB}\right|  = \sqrt{{\left( 4 - 1\right) }^{2} + {\left( -3 - 1\right) }^{2}} = 5$
        \begin{enumerate}
            \item 当 $t = \dfrac{\left| AB\right| }{2} = \dfrac{5}{2}$ ,此时易得符合题意的直线 $l$ 为线段 ${AB}$ 的垂直平分线 ${6x} - {8y} - {23} = 0$以及与直线 ${AB}$ 平行的两条直线 ${8x} + {6y} + {11} = 0$ 和 ${8x} + {6y} - {39} = 0$
            \item 当 $t < \dfrac{\left| AB\right| }{2} = \dfrac{5}{2}$ 时,有 4 条直线 $l$ 会使得点 $A\left( {1,1}\right)$ 和 $B\left( {4, - 3}\right)$ 到它们的距离相等,注意到 $l$ 不过原点,所以当其中一条直线过原点时, 会作为增根被舍去.
            
            设点 $A$ 到 $l$ 的距离为 $d$
            \begin{dingautolist}{172}
                \item 作为增根被舍去的直线 $l$ ,过原点和 $A,B$ 的中点 $M\left( {\dfrac{5}{2}, - 1}\right)$ ,其方程为 ${2x} + {5y} = 0$ ,此时$t = d = \dfrac{7}{\sqrt{29}} < \dfrac{5}{2}$ ,符合
                \item 作为增根被舍去的直线 $l$ ,过原点且与 ${AB}$ 平行,其方程为 ${4x} + {3y} = 0$ ,此时 $t = d = \dfrac{l}{s} < \dfrac{5}{2}$ ,符合
            \end{dingautolist}
        \end{enumerate}
        综上,满足题意的实数 $t$ 为 $\dfrac{5}{2}$ 或 $\dfrac{7}{5}$ 或 $\dfrac{7\sqrt{29}}{29}$ .
    \end{jieda}
    
    \item 在内切圆圆心为 $M$ 的 $\triangle{ABC}$ 中, ${AB} = 3,{BC} = 4,{AC} = 5$ ,在平面 ${ABC}$ 内,过点 $M$ 作动直线 $l$ ,现将 $\triangle{ABC}$ 沿动直线 $l$ 翻折,使翻折后的点 $C$ 在平面 ${ABM}$ 上的射影 $E$ 落在直线 ${AB}$ 上,点 $C$ 在直线 $l$ 上的射影为 $F$ ,则 $\dfrac{\left| EF\right| }{\left| CF\right| }$ 的最小值为\blkda{$8\sqrt{10} - {25}$}.
    \begin{jieda}
        过 $M$ 作 $\bigtriangleup  {ABC}$ 的三边的垂线,设 $\odot  M$ 的半径为 $r$，则 $3 - r + 4 - r = 5$ ,即 $r = 1$ .以 ${BA},{BC}$ 为坐标轴建立平面直角坐标系,如图所示:
        \begin{center}
            \includegraphics[width=0.5\linewidth]{bodyxb//src/zhixianduichen4.png}
        \end{center}    
        则 $M\left( {1,1}\right) ,C\left( {0,4}\right)$ .
        $\because C$ 在平面 ${ABM}$ 的射影在直线 ${AB}$ 上,故直线 $l$ 必存在斜率.
        过 $C$ 作 ${CF} \bot  l$ ,垂足为 $F$ ,交直线 ${AB}$ 于 $E$ .
        设直线 $l$ 的方程为: $y = k\left( {x - 1}\right)  + 1$ ,则 $\left| {CF}\right|  = \dfrac{\left| k + 3\right| }{\sqrt{{k}^{2} + 1}}$ ,
        又直线 ${CE}$ 的方程为: $y =  - \dfrac{1}{k}x + 4$ ,
        $\therefore E\left( {{4k},0}\right) ,\therefore {CE} = \sqrt{{16}{k}^{2} + {16}} = 4\sqrt{{k}^{2} + 1}$ ,
        故而 $\left| {EF}\right|  = 4\sqrt{{k}^{2} + 1} - \dfrac{\left| k + 3\right| }{\sqrt{{k}^{2} + 1}}$ .
        $\therefore \dfrac{\left| EF\right| }{\left| CF\right| } = \dfrac{4\left( {{k}^{2} + 1}\right) }{\left| k + 3\right| } - 1$ .
        \begin{dingautolist}{172}
            \item 若 $k >  - 3$ ,则 $\dfrac{4\left( {{k}^{2} + 1}\right) }{\left| k + 3\right| } - 1 = 4\left( {k + 3}\right)  + \dfrac{40}{k + 3} - {258}\sqrt{10} - {25}$ .
            当且仅当 $4\left( {k + 3}\right)  = \dfrac{40}{k + 3}$ 即 $k = \sqrt{10} - 3$ 时取等号.
            \item 若 $k <  - 3$ ,则 $\dfrac{4\left( {{k}^{2} + 1}\right) }{\left| k + 3\right| } - 1 = 4\left( {-k - 3}\right)  + \dfrac{40}{-k - 3} + {238}\sqrt{10} + {23}$ .
            当且仅当 $4\left( {-k - 3}\right)  = \dfrac{40}{-k - 3}$ 即 $k =  - \sqrt{10} - 3$ 时取等号.
        \end{dingautolist}
        综上可知 $\dfrac{\left| EF\right| }{\left| CF\right| }$ 的最小值为 $8\sqrt{10} - {25}$
    \end{jieda}
\end{enumerate}

\clearpage
\section*{直线综合练习 1}
\begin{enumerate}
    \item 若直线 ${l}_{1}$ 与直线 ${l}_{2} : {3x} - {5y} + 4 = 0$ 平行,且它在两坐标轴上的截距之和为 4,则直线 ${l}_{1}$ 的方程为\blkda[3]{${3x} - {5y} - {30} = 0$}.

    \item 若直线 ${l}_{1}$ 与直线 ${l}_{2} : {3x} + {4y} - 5 = 0$ 垂直,且直线 ${l}_{1}$ 与两坐标轴所构成的三角形面积为 12，则直线 ${l}_{1}$ 的方程为\blkda[7]{ ${4x} - {3y} + {12}\sqrt{2} = 0$ 或 ${4x} - {3y} - {12}\sqrt{2} = 0$}.

    \item 直线 ${4x} - {3y} - 8 = 0$ 的倾斜角的角平分线所在的直线的方程是\blkda[3]{$x - {2y} - 2 = 0$}.

    \item 已知三条直线 ${l}_{1} : {2x} + 1 = 0$ ,直线 ${l}_{2} : {mx} + y = 0$ ,直线 ${l}_{3} : x + {my} - 1 = 0$ . 若这三条直线不能围成三角形,则实数 $m$ 所有可能的值为\blkda[3]{$0, \pm  1, \pm  \sqrt{3}$}.

    \item 已知两条直线 ${l}_{1} : x + y\sin \theta  - 1 = 0$ 和直线 ${l}_{2} : {2x}\sin \theta  + y + 1 = 0$ . 当 $\theta  =$ \blkda[4]{$\dfrac{k}{2}\pi  + \dfrac{\pi }{4}\left( {k \in  \mathbf{Z}}\right) $}时, ${l}_{1}//{l}_{2}$ ；当 $\theta  =$ \blkda[3]{ ${k\pi }\left( {k \in  \mathbf{Z}}\right)$}时, ${l}_{1}\bot {l}_{2}$ .

    \item 点 $P\left( {-1,3}\right)$ 到直线 $l : y = k\left( {x - 2}\right)$ 的距离的最大值等于\blkda{$3\sqrt{2}$}.

    \item 直线 ${3x} - {2y} + m = 0$ 和直线 ${6x} - {4y} + 5 = 0$ 的位置关系是\dotfill{(\kh{B})}
    \xx{平行}{平行或重合}{重合}{既不平行也不重合}


    \item 已知直线 ${l}_{1} : y = x \cdot  \sin \alpha$ 和直线 ${l}_{2} : y = {2x} + c$ ,则直线 ${l}_{1}$ 与 ${l}_{2}$\dotfill{(\kh{D})}
    \xx{通过平移可以重合}{不可能垂直}{可能与 $x$ 轴围成等腰直角三角形}{通过绕 ${l}_{1}$ 上某一点旋转可以重合}
    
    \item 已知直线 ${l}_{1} : x + y\sqrt{1 - \cos a} + b = 0$ ,直线 ${l}_{2} : x\sin a + y\sqrt{1 + \cos a} - a = 0$ . 若 $a \in  \left( {\pi ,\dfrac{3\pi }{2}}\right)$ , 则直线 ${l}_{1}$ 与直线 ${l}_{2}$ 的位置关系是\dotfill{(\kh{B})}
    \xx{平行}{垂直}{平行或重合}{相交但不一定垂直}
    
    \item 若直线 ${l}_{1} : x + {2y} - 3 = 0$ 的倾斜角为 $\theta$ ,且直线 ${l}_{1}$ 与直线 ${l}_{2}$ 关于直线 ${l}_{3} : x = a\left( {a \in  \mathbf{R}}\right)$ 对称,则直线 ${l}_{2}$ 的倾斜角是\dotfill{(\kh{C})}
    \xx{$\dfrac{\pi }{2} - \theta$}{$\theta  - \dfrac{\pi }{2}$}{$\pi  - \theta$}{${2\pi } - \theta$}
    
    \item 已知常数 $\theta  \in  \lbrack  - \pi ,0)$ ,试用 $\theta$ 表示经过 $P\left( {0,0}\right) ,Q\left( {\sin \theta ,\cos \theta }\right)$ 两点的直线 $l$ 的倾斜角\blkda[4]{$\begin{cases}
    -\dfrac{\pi}{2} - \theta, & \theta \in \left[-\pi, -\dfrac{\pi}{2}\right] \\
    \dfrac{\pi}{2} - \theta, & \theta \in \left(-\dfrac{\pi}{2}, 0\right)
    \end{cases}$}.

    \item 已知 $A\left( {-1, - 5}\right) ,B\left( {3, - 2}\right)$ ,直线 $l$ 倾斜角是直线 ${AB}$ 倾斜角的 2 倍,则 $l$ 的斜率\blkda{$\dfrac{24}{7}$}.

    \item 已知直线 $l$ 被直线 $y = {4x} + 6$ 与 $y = {3x} - 7$ 截得的线段中点恰为原点,则 $l$ 的倾斜角\blkda{$arctan(10)$}.

    \item 已知直线 ${l}_{1} : {mx} + {8y} + n = 0$ 与直线 ${l}_{2} : {2x} + {my} - 1 = 0$ 互相平行,垂直于直线 ${l}_{1}$ 的直线 ${l}_{3}$ 过点$(m, n)$,且它被直线 ${l}_{1}$ 与直线 ${l}_{2}$ 所截得的线段长为 $\sqrt{5}$ ,求直线 ${l}_{3}$ 的方程.
    \begin{jieda}
        ${2x} - y + {10} = 0$ 或 ${2x} - y - {30} = 0$ 或 ${2x} + y - {14} = 0$ 或 ${2x} + y + {26} = 0$
    \end{jieda}

    \item 求证: 点 $P\left( {-2,2}\right)$ 到直线 $\left( {2 + k}\right) x - \left( {1 + k}\right) y - 2\left( {3 + {2k}}\right)  = 0$ 的距离不大于 $4\sqrt{2}$ .
    \begin{jieda}
        直线 $\left( {2 + k}\right) x - \left( {1 + k}\right) y - 2\left( {3 + {2k}}\right)  = 0$ 过定点 ${P}_{1}\left( {2, - 2}\right)$，可证得点 $P$ 到该直线的距离不大于 $\left| {P{P}_{1}}\right|  = 4\sqrt{2}$
    \end{jieda}

    \item 对于平面上任意一点 $P$ ,当点 $Q$ 在线段 ${AB}$ 上运动时,称 $\left| {PQ}\right|$ 的最小值为点 $P$ 与线段 ${AB}$ 的距离,记作 $d\left( {P,{AB}}\right)$ .
    \begin{enumerate}
        \item 已知定点 $A\left( {1, - 2}\right)$ 和 $B\left( {4,1}\right)$ ,动点 $P$ 在 $x$ 轴上运动,写出点 $P\left( {t,0}\right)$ 到线段 ${AB}$ 的距离 $h$ 关于 $t$ 的函数关系式
        \item 设 $l$ 是长为 2 的线段,求点的集合 $D = \{ P \mid  d\left( {P,l}\right)  \leq  1\}$ 所表示的图形面积
    \end{enumerate}
    \begin{jieda}
        \begin{enumerate}
            \item $h = \left\{  \begin{array}{ll} \sqrt{{t}^{2} - {2t} + 5}, & t \leq   - 1 \\  \dfrac{\left| {t - 3}\right|}{\sqrt{2}} , &  - 1 < t < 5 \\  \sqrt{{t}^{2} - {8t} + {17}}, & t \geq  5 \end{array}\right.$
            \item $S = \pi  + 4$
        \end{enumerate}
    \end{jieda}
    \item 如图,将边长为 $a$ 的正方形纸片 ${ABCD}$ 沿 ${MN}$ 折叠,使点 $B$ 落在 ${AD}$ 上的点 ${B}_{1}$ 处,点 $C$ 落在点 ${C}_{1}$ 处,设折起的四边形 ${MBCN}$ 的面积为 $S.S$ 是否存在最小值? 若存在,求出它的最小值,并建立平面直角坐标系,求出此时折痕 ${MN}$ 所在的直线方程; 若不存在, 说明理由.
    
    \begin{center}
    \includegraphics[width=0.2\textwidth]{bodyxb/src/bo_d2rdghbef24c73b28d50_2_1117_1052_290_251_0.jpg}
    \end{center}
    \begin{jieda}
        $S$ 存在最小值,最小值为 $\frac{3}{8}{a}^{2},{MN}$ 所在的直线方程为 ${4x} + {8y} - {5a} = 0$
    \end{jieda}

\end{enumerate}